\documentclass[11pt,a4paper]{article}
\usepackage[T1]{fontenc}
\usepackage[utf8]{inputenc}
\usepackage{amsmath,amsthm,mathtools}
\usepackage{newpxtext,newpxmath}
\usepackage[margin=27mm,headheight=15pt]{geometry}
\usepackage{microtype}
\usepackage{booktabs,tabularx,array}
\usepackage[table]{xcolor}
\usepackage{enumitem}
\usepackage{aliascnt}
\usepackage{fancyhdr}
\usepackage{tikz}
\usetikzlibrary{arrows.meta,positioning}
\usepackage[most]{tcolorbox}
\usepackage{listings}
\usepackage{hyperref}
\usepackage[nameinlink,noabbrev]{cleveref}
\definecolor{ink}{HTML}{213F59}
\definecolor{soft}{HTML}{EFF4F7}
\definecolor{muted}{HTML}{52616C}
\hypersetup{colorlinks=true,linkcolor=ink,citecolor=ink,urlcolor=ink,
 pdftitle={Expanding Polynomial Dynamics over Surreal and Surcomplex Fields},
 pdfauthor={Research draft prepared with ChatGPT},
 pdfsubject={Exact symbolic fibers, order-unit scales, Hahn fields, and polynomial iteration}}
\pagestyle{fancy}
\fancyhf{}
\fancyhead[L]{\small\color{muted}Expanding polynomial dynamics}
\fancyhead[R]{\small\color{muted}Hahn and surcomplex fields}
\fancyfoot[C]{\thepage}
\renewcommand{\headrulewidth}{0.3pt}
\setlength{\parskip}{4pt}
\setlength{\parindent}{1.1em}
\setlength{\emergencystretch}{2em}
\setcounter{tocdepth}{2}
\numberwithin{equation}{section}
\newtheorem{theorem}{Theorem}[section]
\theoremstyle{plain}
\newaliascnt{proposition}{theorem}
\newtheorem{proposition}[proposition]{Proposition}
\aliascntresetthe{proposition}
\crefname{proposition}{proposition}{propositions}
\Crefname{proposition}{Proposition}{Propositions}
\newaliascnt{lemma}{theorem}
\newtheorem{lemma}[lemma]{Lemma}
\aliascntresetthe{lemma}
\crefname{lemma}{lemma}{lemmas}
\Crefname{lemma}{Lemma}{Lemmas}
\newaliascnt{corollary}{theorem}
\newtheorem{corollary}[corollary]{Corollary}
\aliascntresetthe{corollary}
\crefname{corollary}{corollary}{corollaries}
\Crefname{corollary}{Corollary}{Corollaries}
\theoremstyle{definition}
\newaliascnt{definition}{theorem}
\newtheorem{definition}[definition]{Definition}
\aliascntresetthe{definition}
\crefname{definition}{definition}{definitions}
\Crefname{definition}{Definition}{Definitions}
\newaliascnt{example}{theorem}
\newtheorem{example}[example]{Example}
\aliascntresetthe{example}
\crefname{example}{example}{examples}
\Crefname{example}{Example}{Examples}
\newaliascnt{question}{theorem}
\newtheorem{question}[question]{Question}
\aliascntresetthe{question}
\crefname{question}{question}{questions}
\Crefname{question}{Question}{Questions}
\theoremstyle{remark}
\newaliascnt{remark}{theorem}
\newtheorem{remark}[remark]{Remark}
\aliascntresetthe{remark}
\crefname{remark}{remark}{remarks}
\Crefname{remark}{Remark}{Remarks}

\newcommand{\No}{\mathbf{No}}
\newcommand{\SC}{\No[\mathrm i]}
\newcommand{\R}{\mathbb R}
\newcommand{\C}{\mathbb C}
\newcommand{\Q}{\mathbb Q}
\newcommand{\Z}{\mathbb Z}
\newcommand{\N}{\mathbb N}
\newcommand{\OO}{\mathcal O}
\newcommand{\mm}{\mathfrak m}
\newcommand{\Ik}{I_\kappa}
\newcommand{\BI}{\mathcal B_{\mathrm{int}}}
\newcommand{\BV}{\mathcal B_{\mathrm{val}}}
\newcommand{\CS}{\mathcal C}
\newcommand{\Sig}{\Sigma_d}
\newcommand{\sh}{\sigma}
\newcommand{\supp}{\operatorname{supp}}
\newcommand{\st}{\operatorname{st}}
\newcommand{\Fix}{\operatorname{Fix}}
\newcommand{\Per}{\operatorname{Per}}
\newcommand{\PrePer}{\operatorname{PrePer}}
\newcommand{\ord}{\operatorname{ord}}
\newcommand{\id}{\operatorname{id}}
\newcommand{\ev}{\operatorname{ev}}
\newcommand{\red}{\operatorname{red}}
\newcommand{\coeff}[2]{[t^{#1}]#2}
\newcommand{\smallbox}[2]{\begin{tcolorbox}[colback=soft,colframe=ink,
 boxrule=0.4pt,arc=1mm,title={#1},fonttitle=\bfseries]#2\end{tcolorbox}}
\lstset{basicstyle=\small\ttfamily,breaklines=true,columns=fullflexible,
 frame=single,rulecolor=\color{muted},showstringspaces=false}

\begin{document}
\hypersetup{pageanchor=false}
\begin{titlepage}
\vspace*{4mm}
{\color{ink}\large RESEARCH REPORT\par}
\vspace{4mm}
{\color{ink}\huge\bfseries Expanding Polynomial Dynamics\par
 over Surreal and Surcomplex Fields\par}
\vspace{6mm}
{\Large Exact symbolic fibers and the order-unit obstruction\par}
\vspace{6mm}
{\large Prepared for Vladimir Reshetnikov\par}
\vspace{2mm}
{\color{muted}22 September 2026\par}
\vspace{2mm}
{\small\color{muted}A single-source report of the collection,
\texttt{docs/surcomplex/expanding-polynomial-dynamics}\par}
\vspace{8mm}
\begin{abstract}
Let $K=k((t^\Gamma))$, with $k=\R$ or $\C$ and $\Gamma$ an arbitrary
nonzero set-sized ordered abelian group. We study
$F=q^{-1}P$, where $v(q)=\kappa>0$,
$\deg P=\deg\red P=d\ge2$, and the reduction of
$P\in\OO[X]$ splits into distinct roots in $k$. We construct, by finite-parameter
formal algebra and strong Hahn evaluation, a canonical point $x_s$ for
every one-sided $d$-symbol itinerary. The entire integral nonescaping set
is the disjoint union $\bigsqcup_s(x_s+I_\kappa)$, where
$I_\kappa=\{h:v(h)>n\kappa\text{ for every }n\in\N\}$.
Different itineraries first separating at time $m$ have distance valuation
exactly $m\kappa$.

This gives an exact distinction between a Cantor repeller and a system
with open infinitesimal fibers. The distinction is whether the expansion
scale $\kappa$ is an order unit of $\Gamma$, not valuation rank alone.
When $\kappa$ is not an order unit the
iterate family is uniformly equicontinuous on the integral nonescaping
set, and every compact forward-invariant subset is finite, although the
scale quotient is the full shift. Every preperiodic point is a canonical
center, and all preperiodicity polynomials split simply in the original
Hahn field. Scalar extension can therefore create entire itinerary
fibers without creating any preperiodic points. We also prove uniform
support-controlled deformation estimates, an exact formula for the
valuation-bounded orbit locus, and a polynomial equicontinuity theorem
when the value group has no order unit. The example
$F(x)=\omega(x^2-1)$ makes all the phenomena explicit inside $\No$ and
$\No[\mathrm i]$: its centers are real surreal numbers, and the fibers in
$\No$ are obtained by restricting to real infinitesimal errors. At the
level of the whole class $\No[\mathrm i]$, the discreteness and
compactness clauses are special cases of the collection's theorem that
every set is closed and discrete in the fine topology; the substance of
those clauses lies in set-sized Hahn workspaces.
\end{abstract}
\vfill
\smallbox{Status and scope}{This is an unrefereed, AI-assisted research
draft with written proofs and exact finite regression tests, placed in the
collection from one delivered manuscript (\cref{epd:subsec:provenance}).
The fiber, compact-subsystem,
and scalar-extension/deformation results are proposed as new results;
historical priority is not certified. No named published open conjecture
is claimed solved. The scale subgroup and error-ideal lemma are formalized in Lean;
the dynamical theorems remain pending. See \cref{epd:subsec:formalization}.}
\end{titlepage}
\pagenumbering{roman}
\hypersetup{pageanchor=true}
\begingroup
\small
\setlength{\parskip}{0pt}
\tableofcontents
\endgroup
\clearpage
\pagenumbering{arabic}

\section{The problem and the contribution}
\label{epd:sec:intro}
A familiar construction in non-Archimedean dynamics takes finitely many
inverse branches, each contracting by a fixed scale, and labels the
points that never escape by infinite words. Over a rank-one valued field,
the contraction scales tend to zero, and a word determines a single
point. Classical quadratic Cantor Julia sets are an established example;
see Benedetto's lecture notes, Example~4.39 \cite{Benedetto}.

The same conclusion need not survive in a larger Hahn field, and it does
not survive unchanged in the full surreal or surcomplex field. If
$\kappa>0$ is the valuation gained by one inverse branch, then the gains
$n\kappa$ may remain below a larger valuation $\eta$ for every ordinary
integer $n$. Successive choices resolve a point to every $q$-order, but
may never resolve a perturbation of valuation $\eta$. There are two
separate questions: whether every word is realized at all, and exactly
how many points realize it. A topological contraction argument alone
answers neither question at arbitrary rank.

We answer both questions for a substantial polynomial family, including
arbitrary infinitesimal coefficient perturbations of degree at most $d$
of a degree-$d$ polynomial with simple split reduction. The main formula is
\begin{equation}\label{epd:eq:intro-fibers}
 \BI(F)=\coprod_{s\in\Sig}(x_s+\Ik),
 \qquad
 \Ik=\bigcap_{n\geq0}\{h\in K:v(h)>n\kappa\}.
\end{equation}
The right-hand side describes actual points, not only formal labels or
an unspecified intersection of disks. The centers $x_s$ are constructed
with one support certificate valid for all itineraries. Their
construction does not presume topological convergence of a sequence of
Hahn elements.

The result has several consequences which are useful precisely when the
rank-one picture breaks down. The preperiodic points remain rigid while
the fibers grow. The native topology can make all compact invariant
subsystems finite while a coarser scale quotient has entropy $\log d$.
Perturbations beyond every power of $q$ preserve the whole itinerary
partition, yet move the actual periodic points. Finally, leaving the
integral disk and having a valuation-unbounded orbit are different
properties; we calculate exactly when they agree.

\subsection{Relation to the repository at the pinned revision}
\label{epd:subsec:repo-pin}
The repository used is \texttt{VladimirReshetnikov/Surreal}, at revision
\begin{center}
\small\texttt{a3124af79f66b8b9c196d76b4cbc5ac3938907c4}.
\end{center}
Its root and documentation inventories
were read, along with the description of
\nolinkurl{docs/surcomplex/dynamics-and-normal-forms}. That report concerns
local normal forms, return thresholds, small divisors, and associated
formal/support questions. The present article studies a global integral
nonescaping set and its infinite symbolic fibers instead \cite{Repo,RepoDynamics}.

This was a targeted inventory audit, not a line-by-line review of every
long manuscript or Lean file. The project already develops Hahn support
machinery and distinctions among analytic categories. Those ideas are
not presented here as discoveries. We use a classical Hahn-field
workspace, state every support hypothesis, and give the polynomial and
topological arguments in full.
The relation to the reports of the collection as it now stands, including
reports the manuscript did not consult, is recorded in
\cref{epd:subsec:collection}.

\subsection{What is, and is not, claimed to be new}
The positive-support lemma, the formal implicit-function construction,
valuation coarsening, and the one-sided full shift are standard tools.
The rank-one Cantor construction and the elementary periodic-word counts
are also standard. Our proposed new contribution is the combined
arbitrary-rank theorem: exact affine itinerary fibers; the sharp
order-unit/topology dichotomy; finiteness of compact forward-invariant
subsets when the scale is not an order unit; and the scalar-extension and
support-controlled deformation consequences. (Placement note: this claim
concerns set-sized Hahn workspaces. For the whole class $\No[\mathrm i]$
with its fine topology, the discreteness and compactness clauses are
special cases of a theorem already in the collection; see
\cref{epd:subsec:collection} and \cref{epd:rem:discrete-special}.)

A targeted search through public sources on Hahn fields,
non-Archimedean polynomial dynamics, higher-rank valuations, and surreal
analysis did not locate these statements in this form. That is not a
proof of historical originality. The article therefore supplies proofs,
not a claim of independently certified priority. It does not purport to
resolve the unrelated birthday, automorphism, or differential-equation
questions in the repository.

\smallbox{Three distinctions that cannot be suppressed}{
An element of negative derivative valuation is called
\emph{valuation-expanding} here; this does not automatically make its
periodic point repelling for a real-valued absolute value. An
\emph{integral nonescaping} orbit stays in $\OO$; a
\emph{valuation-bounded} orbit merely has some lower valuation bound.
Finally, strong Hahn summation is an algebraic support operation, not
necessarily a limit in the full valuation topology.}

\subsection{Provenance of this report}
\label{epd:subsec:provenance}
This report is built from \textbf{one} delivered manuscript, the fifth item
of the eighteenth batch of manuscripts received by the collection (archive
name \texttt{surreal\_symbolic\_dynamics}). It was delivered on
22 September 2026 as a 24-page research draft; its PDF metadata describe it
as a research draft prepared with ChatGPT. The delivered title was
\emph{Expanding Polynomial Dynamics over Surreal and Surcomplex Fields:
Exact symbolic fibers and the cofinality obstruction}. Its repository pin is
the revision \texttt{a3124af} displayed in
\cref{epd:subsec:repo-pin}. The cross-references added at placement were
checked against the collection at commit \texttt{e9a9650}.

This is not a merge. There was no second source, and nothing was selected
out of a larger body of work: every theorem, proof, example, question, and
limitation of the manuscript is kept. The placement changed the following,
and nothing else.
\begin{enumerate}[label=\textup{(\arabic*)},leftmargin=2em]
\item Every label carries the prefix \texttt{epd:}.
\item The manuscript called an order unit a \emph{cofinal scale} and spoke of
the \emph{cofinal} and \emph{noncofinal} cases. That word is reserved in
the collection for another invariant, so the property is called an order unit
throughout. The subtitle and the name of \cref{epd:thm:topology} change
accordingly. The dictionary is stated once, in \cref{epd:rem:one-name}.
\item Cross-notes to neighbouring reports were added
(\cref{epd:subsec:collection,epd:rem:discrete-special}), together with a
pointer to the collection's Lean proof of the positive-support lemma
(\cref{epd:subsec:formalization}).
\item The reproduction instructions were adapted to the placed file layout.
The verification program overwrites its own recorded output by default,
and \cref{epd:subsec:reproduce} says how to avoid that.
\item The non-claims scattered through the manuscript are collected in
Appendix~\ref{epd:app:nonclaims}.
\end{enumerate}
At placement, no hypothesis, statement, or proof was altered beyond the
change of word in~(2). A subsequent proof review made degree preservation
explicit in the summaries, added its counterexample below
\eqref{epd:eq:polynomial}, expanded the noncompactness and formal-substitution
arguments, and updated Lean coverage. The delivered code and audit remain
unchanged. The report is filed under \texttt{surcomplex} although the archive
name says \emph{surreal}. Its theorems hold uniformly for $k=\R$ and
$k=\C$, the explicit specialization (\cref{epd:thm:surreal}) is stated in
$\No[\mathrm i]$, and its valuation-theoretic neighbours are surcomplex
reports. The real case is kept explicit throughout. For $k=\R$ the
hypothesis \eqref{epd:eq:polynomial} asks for real, distinct residue
roots. The centers of the surreal example are real, and its fibers in
$\No$ are obtained by restricting to real errors.

\subsection{Relation to other reports in this collection}
\label{epd:subsec:collection}
The manuscript consulted only the repository inventories and the
description of the dynamics report. The following notes place it among the
reports the collection now contains. They were written at placement. No
result is transferred into this report from any of them, and none of
their results is strengthened here.

\paragraph{\emph{Surcomplex Dynamics: Linearization Thresholds, Periods,
and Normal Forms}} (\texttt{docs/\allowbreak surcomplex/\allowbreak
dynamics-and-\allowbreak normal-forms/}). The two reports are about different
problems. That report linearizes \emph{germs} at a fixed point whose
multiplier is an ordinary complex number of modulus one, with small
divisors, return thresholds, and coefficient-function categories. Its
limitations section excludes ``global Julia or Fatou theory on the proper
class $\No[\mathrm i]$'' and, from the scope of its exact finite-return
linearization theorem (\texttt{dyn:thm:exact}), ``attracting or repelling
multipliers'' (items N51 and N18 of its non-claims list when this note was
written). This
report answers \emph{neither} exclusion. It studies the global integral
nonescaping set of a polynomial whose periodic points have
valuation-expanding multipliers. It constructs no Julia or Fatou theory:
$J_{\mathrm{aff},v}$ of \cref{epd:def:affine-julia} is a deliberately
named affine equicontinuity set, and \cref{epd:cor:empty-affine} says
nothing about spherical or Berkovich Julia sets. By \cref{epd:rem:expansion},
``valuation-expanding'' is not ``repelling''. It linearizes nothing at any
periodic point. Several letters and words mean different things in the
two reports, and each keeps its own meaning:
\begin{itemize}[leftmargin=1.4em]
\item $\sigma$ is the one-sided shift here. There $\sigma(\lambda)$ is the
drifted exponential small-divisor rate.
\item $\kappa=v(q)>0$ is the valuation gained by an inverse branch here.
There it is the generic divisor rate of \texttt{dyn:rem:oneinvariant}, and
also appears as $\kappa_\Upsilon$.
\item $q$ is a Hahn element of valuation $\kappa$ here. There it is the order
in the resonance tag $\mathsf{RES}_q$, and $q_j$ are continued-fraction
denominators.
\item A \emph{multiplier} here is $(F^p)'(x)\in K$ at a periodic point, a Hahn
element of valuation $-p\kappa<0$ (\cref{epd:prop:multiplier}). There
it is an ordinary unit complex number.
\item \emph{Julia} occurs here only in the subscripted affine set
$J_{\mathrm{aff},v}$. There it refers to Julia's functional equation
(\texttt{dyn:cor:julia}).
\end{itemize}
The quadratic critical scale of that report
(\texttt{dyn:prop:quadratic-scale}) rescales $\lambda z+t^\eta z^2$ at a
unit multiplier, as a boundary of a linearization domain. It is unrelated
to the expanding scale $q^{-1}$ used here.

\paragraph{\emph{Surcomplex Analysis}} (\texttt{docs/\allowbreak
surcomplex/\allowbreak analysis/}). Its fine topology
(\texttt{a:def:fine}) is the topology that the manuscript calls the full
native, or order (fine), topology of $\No$ and $\No[\mathrm i]$ in
\cref{epd:sec:surreal}. Its \texttt{a:thm:discrete} proves that every
\emph{set} $A\subseteq\No[\mathrm i]$ is closed and discrete in that
topology, with one separating radius for all pairs of points of~$A$. The
surreal-level discreteness and compactness statements of this report are
therefore special cases: see \cref{epd:rem:discrete-special}. The
substance of \cref{epd:thm:topology,epd:thm:compact} lies in set-sized
workspaces $k((t^\Gamma))$ with their intrinsic valuation topologies. There
a set need not be discrete, and the finiteness of compact invariant
subsets is not automatic. The manuscript's remark comparing the fine topology
with a workspace topology (\cref{epd:rem:fine-vs-workspace}) is an instance
of \texttt{a:rem:notvaluation}. The uniform equicontinuity statements over
whole fine balls, which are proper classes, are not consequences of
\texttt{a:thm:discrete}.

\paragraph{\emph{Cofinality, Factorization and Scalar Extension for Entire
Hahn Functions at Arbitrary Rank}} (\texttt{docs/\allowbreak
surcomplex/\allowbreak entire-functions-\allowbreak at-arbitrary-rank/}). An order unit
in \cref{epd:def:orderunit} is exactly an order unit in the sense of
\texttt{ent:def:order-unit}. That report reserves the word
\emph{cofinality} for $\operatorname{cf}(\Gamma)$
(\texttt{ent:rem:one-name}), which is why this report does not call an order
unit a cofinal scale. The maximal proper convex subgroup $H$ of
\cref{epd:subsec:coarsened} is the kernel of the order-preserving
homomorphism $\pi:\Gamma\to\R$ of \texttt{ent:lem:coarsening}. That lemma is
stated under that report's standing divisibility hypothesis. The facts about
$H$ used here are proved directly in \cref{epd:subsec:coarsened} and need no
divisibility. The countable group without an order unit in
\cref{epd:ex:no-order-unit} is, as an ordered group, that report's
$\Gamma_\infty=\bigoplus_{j\geq0}\Q\,\omega^j$
(\texttt{ent:subsec:gamma-infinity}). Its \texttt{ent:\allowbreak cor:cofinality} is a
dichotomy in $\operatorname{cf}(\Gamma)$ for entire functions. It is not
the order-unit dichotomy of \cref{epd:thm:topology}.

\paragraph{\emph{Hahn--Tate Uniformization at Arbitrary Valuation Rank}}
(\texttt{docs/\allowbreak surcomplex/\allowbreak hahn-tate-\allowbreak uniformization/}).
The convex subgroup $H_\kappa$ of \cref{epd:def:orderunit} is the
subgroup $H_\alpha$ of \texttt{tate:eq:H}, for $\alpha=\kappa$: the smallest
convex subgroup containing the scale. There it bounds the domain of
the Tate parametrization. Here it bounds the error ideal $\Ik$. No
result is shared.

\paragraph{\emph{Rank-One Non-Archimedean Analytic Geometry inside the
Surcomplex Numbers}} (\texttt{docs/\allowbreak surcomplex/\allowbreak
rank-one-\allowbreak berkovich/}). When $\kappa$ is an order unit,
\cref{epd:thm:topology}(a) has the same symbolic Cantor model as the
classical rank-one construction cited in
\cite[Example~4.39]{Benedetto}. An order unit does not make $\Gamma$
rank one: $\Q\oplus_{\mathrm{lex}}\Q$ with $\kappa=(1,0)$ is already
a counterexample. The citation provides the rank-one quadratic comparison;
the proof here uses the exact cylinder valuations even at higher rank.
In the fixed field $\C((t^{\R}))$ of that report every positive element of
the Archimedean group $\R$ is an order unit, so only that case can occur
there. That report contains no polynomial dynamics. No statement about
Berkovich spaces, Julia sets, or Tate algebras is proved or used here, and
no transfer in either direction is claimed.

\section{Workspaces, scales, and support}
\label{epd:sec:workspace}
\subsection{The valued field}
Throughout Sections~\ref{epd:sec:workspace}--\ref{epd:sec:general}, let
\[
 k\in\{\R,\C\},\qquad K=k((t^\Gamma)),
\]
where $\Gamma\ne0$ is a set-sized ordered abelian group. An element is a
series $a=\sum_{\gamma\in S}a_\gamma t^\gamma$ with $S\subseteq\Gamma$
well ordered. All coefficients in $k$ are trivially valued. Set
\[
 v(a)=\min\supp(a)\quad(a\ne0),\qquad v(0)=+\infty,
 \qquad \OO=\{a:v(a)\geq0\},\qquad \mm=\{a:v(a)>0\}.
\]
The residue map $\st:\OO\to k$ extracts the coefficient at $0$.
The usual complex modulus of a coefficient is not the valuation used
in this article. No divisibility assumption on $\Gamma$, and no
algebraic closedness of $K$, is needed for the main theorem.

We use $\N=\{0,1,2,\ldots\}$. The native valuation topology has balls
$\{y:v(y-x)>\gamma\}$, for $\gamma\in\Gamma$, as a neighborhood basis.
The associated additive uniformity has entourages
$\{(x,y):v(x-y)>\gamma\}$. These definitions make sense without a
real-valued metric.

\subsection{An individual scale, not the rank}
\begin{definition}\label{epd:def:orderunit}
A positive $u\in\Gamma$ is an \emph{order unit} if for every
$\gamma\in\Gamma$ there is an
$n\in\N$ with $\gamma\leq nu$. For $\kappa>0$, define
\[
 H_\kappa=\{\gamma\in\Gamma:|\gamma|\leq n\kappa
                    \text{ for some }n\in\N\},\qquad
 \Ik=\{h\in K:v(h)>n\kappa\text{ for every }n\in\N\}.
\]
The convention $v(0)=+\infty$ includes zero in $\Ik$.
When the scale $\kappa=v(q)$ of a map is an order unit we speak of the
\emph{order-unit case}. Otherwise we speak of the \emph{non-order-unit
case}.
\end{definition}

\begin{remark}[One name for one property]\label{epd:rem:one-name}
The delivered manuscript called an order unit a \emph{cofinal scale}. It
called the two cases the \emph{cofinal} and \emph{noncofinal} cases, and
the distinction the \emph{cofinality} of the expansion scale. In this
collection, the report \texttt{entire-functions-\allowbreak at-arbitrary-rank}
(\texttt{ent:rem:one-name}) reserves \emph{cofinality} for
$\operatorname{cf}(\Gamma)$, the least cardinality of a cofinal subset of
$\Gamma$. Its \texttt{ent:def:order-unit} is the notion defined above.
This report therefore uses the following dictionary, and only it.
\begin{center}\small
\begin{tabular}{@{}ll@{}}
\toprule
Delivered manuscript & This report\\\midrule
$\kappa$ is cofinal; a cofinal scale & $\kappa$ is an order unit; an order-unit scale\\
$\kappa$ is noncofinal & $\kappa$ is not an order unit\\
the cofinal / noncofinal case & the order-unit / non-order-unit case\\
cofinality dichotomy (\cref{epd:thm:topology}) & order-unit dichotomy\\
cofinality is lost (\cref{epd:thm:extension}) & the order-unit property is lost\\
\bottomrule
\end{tabular}
\end{center}
The word \emph{cofinal} survives only in its subset sense. A subset
$S$ of $\Gamma$ (or of $\Gamma_{>0}$) is cofinal if every element is at
most some member of $S$. Thus $u>0$ is an
order unit exactly when $\{nu:n\in\N\}$ is cofinal in $\Gamma$, in agreement
with \texttt{docs/NOTATION.md}. Apart from the quotations of the
manuscript in this remark and in \cref{epd:subsec:provenance}, wherever
\emph{cofinality} occurs in this report it means
$\operatorname{cf}(\Gamma)$. It occurs in \cref{epd:ex:no-order-unit} and
in names taken from the entire-functions report. The delivered file \texttt{RESEARCH\_STATUS.md}, kept unchanged, still uses
the manuscript's words.
\end{remark}

\begin{lemma}\label{epd:lem:ideal}
$H_\kappa$ is a convex subgroup of $\Gamma$. The set $\Ik$ is an ideal
of $\OO$, and, for any $q$ with $v(q)=\kappa$, one has $q\Ik=\Ik$.
Moreover
\[
 \Ik=\bigcap_{n\geq0}q^n\OO,
 \qquad
 \Ik=0\quad\Longleftrightarrow\quad \kappa\text{ is an order unit}.
\]
If $\kappa$ is not an order unit, $\Ik$ is a nonzero clopen additive
subgroup of $K$. It is the maximal ideal of the valuation ring for the
coarsened valuation $K^\times\to\Gamma/H_\kappa$.
\end{lemma}
\begin{proof}
Closure and convexity of $H_\kappa$ follow by adding the finite bounds.
The valuation inequality for a sum and the inequality
$v(ah)\geq v(h)$ for $a\in\OO$ prove the ideal assertion. Adding or
subtracting $\kappa$ preserves the condition of exceeding every
$n\kappa$, proving $q\Ik=\Ik$. Requiring $v(h)\geq n\kappa$ for all
$n$ is equivalent to requiring $v(h)>n\kappa$ for all $n$, by increasing
$n$ by one. This proves the intersection formula.

If $\kappa$ is an order unit, no finite valuation can exceed all $n\kappa$.
Otherwise there exists $\eta\in\Gamma$ with $\eta>n\kappa$ for all
$n$, and $t^\eta\in\Ik\setminus\{0\}$. The ball $\{v>\eta\}$ is
contained in $\Ik$, so this subgroup is open, and every open subgroup
of a topological group is closed. Finally,
$v(h)+H_\kappa>0$ means precisely that $v(h)$ exceeds every element of
$H_\kappa$, equivalently every $n\kappa$.
\end{proof}

The ideal is thus not an ad hoc error class. It is the exact error left
unresolved by the convex subgroup generated by the dynamical scale.
It is generally an ideal of a valuation ring, not an ideal of the field
$K$.

\begin{example}\label{epd:ex:ranktwo}
In $\Gamma=\Q\oplus_{\mathrm{lex}}\Q$, the scale $(0,1)$ is not
an order unit, since $(1,0)>n(0,1)$ for every $n$. In the very same group,
$(1,0)$ is an order unit. Consequently ``rank greater than one'' cannot replace
``the chosen scale is not an order unit'' in any theorem below.
\end{example}

\subsection{The summation tool and its exact role}
A family $(a_j)_{j\in J}$ of Hahn series is \emph{strongly summable} if
its support union is well ordered and every exponent occurs in only
finitely many members of the family. Its sum is then defined
coefficientwise.

We use the classical positive-support lemma associated with Neumann's
ordered-series construction \cite{Neumann}: if $S\subseteq\Gamma_{>0}$
is well ordered, the additive monoid $\langle S\rangle$ is well ordered,
and each fixed element of $\Gamma$ has only finitely many
representations as a finite word in $S$, allowing all finite word
lengths. Repetitions and ordering of letters are allowed in this
finiteness assertion. The following useful form spells out the only
infinite evaluation operation needed here.

\begin{lemma}[Finite-parameter Hahn evaluation]\label{epd:lem:evaluation}
Let $a_1,\ldots,a_r\in\mm$. Every formal series
$A\in k[[Z_1,\ldots,Z_r]]$ has a well-defined strong evaluation
$A(a_1,\ldots,a_r)\in\OO$. Evaluation is a $k$-algebra homomorphism.
Every evaluated support is contained in
\[
 M=\left\langle\bigcup_{j=1}^r\supp(a_j)\right\rangle.
\]
Evaluation respects formal substitution by series with zero constant
term. If a formal series belongs to an ideal $(Z_{j_1},\ldots,Z_{j_s})$,
its evaluation belongs to the ideal $(a_{j_1},\ldots,a_{j_s})\OO$.
\end{lemma}
\begin{proof}
The finite support union is well ordered and positive. Expand every
parameter monomial into words in this union, marking a letter by the
parameter from which it came. The positive-support lemma bounds both
the possible words and their lengths at each output exponent. There
are only finitely many choices of a parameter mark at every letter.
Thus the expanded family has well-ordered support and finite
coefficient fibers. Addition, multiplication, and zero-constant formal
substitution may all be regrouped coefficientwise, because at each
exponent only finitely many terms participate.

For the ideal assertion, write the formal series as a finite sum
$\sum_{\ell=1}^s Z_{j_\ell}B_\ell$. Each $B_\ell$ evaluates into
$\OO$ by the first assertion. Apply multiplicativity.
\end{proof}

\begin{remark}\label{epd:rem:no-limits}
The lemma does not assert that partial sums converge in the native
topology. For instance, in the rank-two example with $q=t^{(0,1)}$,
$\sum_{n\geq0}q^n=(1-q)^{-1}$ is a valid Hahn identity, while the
remainders have valuations $(0,N)$ and never enter the neighborhood
$\{v>(1,0)\}$. Algebraic identities must be established before
parameter evaluation; an unproved continuity assertion cannot replace
that step.
\end{remark}

\section{Simple reduction and universal inverse branches}
\label{epd:sec:branches}
Fix $q\in\mm\setminus\{0\}$, write $\kappa=v(q)$, and let
\begin{equation}\label{epd:eq:polynomial}
 P(X)=P_0(X)+\sum_{j=0}^d e_jX^j\in\OO[X],
 \qquad e_j\in\mm,
 \qquad P_0(X)=a_d\prod_{i=1}^d(X-c_i),
\end{equation}
where $d\geq2$, $a_d\in k^\times$, and the $c_i\in k$ are distinct.
Zero perturbation coefficients are permitted. In particular $P$ has
degree $d$ and unit leading coefficient. The equality
$\deg P=\deg P_0$ is essential for the exterior-orbit and global
preperiodicity claims. For example, over $\R((t^\Z))$ take
\[
 q=t,\qquad P(X)=(X^2-1)(1+tX).
\]
Its reduction still has the two simple roots $\pm1$, but $\deg P=3$.
For $F=t^{-1}P$, one has $F(-t^{-1})=0$ and $F(0)=-t^{-1}$.
This gives a two-cycle with a nonintegral point, so the claims below
cannot be extended merely by allowing infinitesimal terms of higher degree.

Primes below always denote
formal differentiation with respect to $X$, holding coefficients fixed;
no differential-field derivation on the surreal coefficients is used. Set
\[
 F(X)=q^{-1}P(X),\qquad
 \BI(F)=\{x\in K:F^n(x)\in\OO\text{ for all }n\in\N\}.
\]
For $x\in\BI(F)$, the identity $P(F^n(x))=qF^{n+1}(x)$ implies
$P_0(\st(F^n(x)))=0$. Thus there is a well-defined itinerary
\[
 \pi(x)=(s_0,s_1,\ldots)\in\Sig:=\{1,\ldots,d\}^{\N},
 \qquad \st(F^n(x))=c_{s_n}.
\]
We write $\sh(s_0,s_1,\ldots)=(s_1,s_2,\ldots)$.

\begin{lemma}[Universal implicit branch]\label{epd:lem:formalbranch}
For each $i$, there exists a unique
$\Psi_i\in k[[E_0,\ldots,E_d,V]]$ such that
\begin{equation}\label{epd:eq:formalbranch}
 \Psi_i(0,0)=c_i,
 \qquad
 P_0(\Psi_i)+\sum_{j=0}^dE_j\Psi_i^j=V.
\end{equation}
Moreover, there is a universal
$A_i\in k[[E_0,\ldots,E_d,U,W]]$ for which
\begin{equation}\label{epd:eq:formal-difference}
 \Psi_i(E,U)-\Psi_i(E,W)=(U-W)A_i(E,U,W),
\end{equation}
where $A_i$ is a unit with constant term $1/P_0'(c_i)$.
The identity may be substituted into another formal power-series ring
whenever all substituted parameters have positive total order.
\end{lemma}
\begin{proof}
Put $\Psi_i=c_i+Y$, with $Y$ of positive total degree. At total degree
$r\geq1$, the only occurrence of the unknown degree-$r$ part of $Y$
is its product with $P_0'(c_i)$. Every other term uses already determined
lower-degree parts, or has a positive-degree parameter as an additional
factor. Since $P_0'(c_i)\ne0$, this determines $Y$ recursively and proves
existence and uniqueness.

Subtract the two equations in \eqref{epd:eq:formalbranch}. The polynomial
divided difference in the two branch values has constant term
$P_0'(c_i)$, hence is a unit. Its inverse gives
\eqref{epd:eq:formal-difference}.
\end{proof}

\begin{proposition}[Exact branch metric]\label{epd:prop:branches}
For every $w\in\OO$ and every $i$, there is a unique $y\in\OO$ with
$\st(y)=c_i$ and $P(y)=qw$. Denote it by $\Phi_i(w)$. For all
$w,z\in\OO$,
\begin{equation}\label{epd:eq:branch-metric}
 v\bigl(\Phi_i(w)-\Phi_i(z)\bigr)=\kappa+v(w-z).
\end{equation}
If $x,y\in\OO$ have the same residue $c_i$, then
\begin{equation}\label{epd:eq:forward-metric}
 v(F(x)-F(y))=v(x-y)-\kappa.
\end{equation}
Equal inputs are interpreted using $+\infty$.
\end{proposition}
\begin{proof}
Evaluate $\Psi_i(E,V)$ at $(e_0,\ldots,e_d,qw)$, all nonzero inputs
having positive valuation. \Cref{epd:lem:evaluation} proves existence and
the defining equation. The residue is $c_i$.

For uniqueness and the exact metric, use the polynomial identity
\[
 P(Y)-P(Z)=(Y-Z)D_P(Y,Z).
\]
When both arguments have residue $c_i$, the residue of $D_P$ is
$P_0'(c_i)\ne0$, so $v(D_P)=0$. Substituting $Y=\Phi_i(w)$ and
$Z=\Phi_i(z)$ gives
$v(Y-Z)=v(q(w-z))$. Substituting arbitrary $x,y$ in the same residue
class gives \eqref{epd:eq:forward-metric}. It also proves that two lifts of
the same $w$ coincide.
\end{proof}

In particular, $F$ maps each of its $d$ integral inverse-branch pieces
bijectively onto $\OO$. Different pieces have points whose differences
have valuation $0$. The branch calculation is algebraic; spherical
completeness or a Banach fixed-point theorem has not been used.

\section{Canonical realization of every itinerary}
\label{epd:sec:centers}
Introduce the finite-parameter ring
\[
 R=k[[Q,E_0,\ldots,E_d]],\qquad
 \mathcal P(X)=P_0(X)+\sum_{j=0}^dE_jX^j.
\]
For $U\in R$, let
$\boldsymbol\Phi_i(U)=\Psi_i(E,QU)$. Substitution is valid since $QU$
has zero constant term. By \cref{epd:lem:formalbranch},
\begin{equation}\label{epd:eq:q-contraction}
 \boldsymbol\Phi_i(U)-\boldsymbol\Phi_i(W)\in Q(U-W)R.
\end{equation}
For $s\in\Sig$, form the depth-$N$ expression
\[
 X_s^{[N]}=
 \boldsymbol\Phi_{s_0}\circ\cdots\circ
 \boldsymbol\Phi_{s_{N-1}}(0),
 \qquad X_s^{[0]}=0.
\]

\begin{theorem}[Universal centers and common support]\label{epd:thm:centers}
There is a formal series $X_s\in R$ for every $s\in\Sig$ with
\begin{equation}\label{epd:eq:formalcenter}
 X_s-X_s^{[N]}\in Q^N R,
 \qquad \mathcal P(X_s)=QX_{\sh s},
 \qquad X_s(0)=c_{s_0}.
\end{equation}
After evaluation $Q=q$, $E_j=e_j$, set $x_s=\ev(X_s)$. Then
\begin{equation}\label{epd:eq:shift-center}
 F(x_s)=x_{\sh s},\qquad \pi(x_s)=s,
\end{equation}
and, simultaneously for all $s$,
\begin{equation}\label{epd:eq:common-support}
 \supp(x_s)\subseteq
 M:=\left\langle\supp(q)\cup\bigcup_{j=0}^d\supp(e_j)\right\rangle.
\end{equation}
Every coefficient of total parameter degree $r$ in $X_s$ depends only
on $s_0,\ldots,s_r$.
\end{theorem}
\begin{proof}
Repeated use of \eqref{epd:eq:q-contraction} shows that for $L\geq N$,
\[
 X_s^{[L]}-X_s^{[N]}\in Q^N R.
\]
The ring $R$ is $Q$-adically complete and separated: grouping a formal
series by its $Q$-degree identifies it with
$k[[E_0,\ldots,E_d]][[Q]]$. For each fixed $Q$-degree, the entire
coefficient series in $E$ stabilizes once $N$ exceeds that degree.
These stabilized coefficients define a unique $X_s$, with the displayed
congruence for every $N$.
This completion takes place in an ordinary finite-variable formal
ring, not in the valuation topology of $K$.

The finite expressions satisfy
$\mathcal P(X_s^{[N+1]})=QX_{\sh s}^{[N]}$. Passing to their formal
$Q$-adic limits proves the middle identity in
\eqref{epd:eq:formalcenter}. The other two assertions follow directly from
the construction. If $N=r+1$, the difference is divisible by $Q^{r+1}$,
so every coefficient of total degree at most $r$ has already stabilized
and only the first $r+1$ symbols have been used.

Evaluate the established formal polynomial identity by
\cref{epd:lem:evaluation}. This proves \eqref{epd:eq:shift-center}. Reduction
at each shifted word gives the itinerary. The same evaluation lemma
gives \eqref{epd:eq:common-support}, with no dependence of $M$ on $s$.
\end{proof}

The formal section is distinguished by the fixed coefficient embedding
$k\hookrightarrow K$, the polynomial reduction $P_0$, and the displayed
universal construction. We do not assert an intrinsic uniqueness among
all possible set-theoretic sections of $\pi$.

\begin{corollary}[Constant coefficients]\label{epd:cor:constant-center}
If $P=P_0$ and $q=t^\kappa$, then each $x_s$ has the form
$\sum_{n\geq0}a_n(s)t^{n\kappa}$ with $a_n(s)\in k$. The coefficient
$a_n(s)$ depends only on the first $n+1$ symbols. The map $s\mapsto x_s$
is unchanged when the exponent group is enlarged.
\end{corollary}
\begin{proof}
With all $E_j=0$, $X_s$ belongs to $k[[Q]]$. Evaluate coefficientwise.
The same formal series is evaluated under every compatible Hahn-field
embedding.
\end{proof}

\begin{remark}[A limit that is not taken]\label{epd:rem:nestedlimit}
For $P=X^2-1$, label roots by signs $\pm1$. The depth-$N$ nested branch
with terminal zero differs from $x_s$ by an element of valuation exactly
$N\kappa$: apply \eqref{epd:eq:branch-metric} $N$ times to the tail $x_{\sh^N s}$
and $0$, whose residues differ. When $\kappa$ is not an order unit, these
approximants do not converge to $x_s$ in the native valuation topology.
Nevertheless $x_s$ is a perfectly well-defined Hahn element by the
preceding theorem.
\end{remark}

\section{The exact symbolic-fiber theorem}
\label{epd:sec:fibers}
\begin{theorem}[Exact fibers and exact separation]\label{epd:thm:fibers}
Under \eqref{epd:eq:polynomial}, the itinerary map is onto and
\begin{equation}\label{epd:eq:fiber-main}
 \pi^{-1}(s)=x_s+\Ik,
 \qquad
 \BI(F)=\coprod_{s\in\Sig}(x_s+\Ik).
\end{equation}
For $h\in\Ik\setminus\{0\}$ and every $n\geq0$,
\begin{equation}\label{epd:eq:fiber-growth}
 v\bigl(F^n(x_s+h)-x_{\sh^n s}\bigr)=v(h)-n\kappa.
\end{equation}
If $s\ne r$ first differ at index $m$, then for every
$x\in\pi^{-1}(s)$ and $y\in\pi^{-1}(r)$,
\begin{equation}\label{epd:eq:first-mismatch}
 v(x-y)=m\kappa.
\end{equation}
Thus a word determines a unique point if and only if $\kappa$ is an
order unit in the ambient value group.
\end{theorem}
\begin{proof}
Surjectivity is already supplied by the centers. Let $h\in\Ik$.
Initially $x_s+h$ and $x_s$ have the same residue. If their $n$th
iterates have difference valuation $v(h)-n\kappa$, that valuation is
positive by the definition of $\Ik$. They therefore have the same
residue at time $n$, and \eqref{epd:eq:forward-metric} proves the next
instance of \eqref{epd:eq:fiber-growth}. Induction proves both that the
perturbed point remains integral and that its itinerary is $s$.

Conversely, let $x$ have itinerary $s$. For every $N$ we may recover
$x$ and $x_s$ by applying the same first $N$ inverse branches to
$F^N(x)$ and $x_{\sh^N s}$, respectively. These two terminal points are
integral, so their difference has nonnegative valuation. Applying
\eqref{epd:eq:branch-metric} $N$ times gives
\[
 v(x-x_s)\geq N\kappa.
\]
This holds for every $N$, and \cref{epd:lem:ideal} implies $x-x_s\in\Ik$.
No limiting argument is involved.

For two different itineraries, their $m$th iterates have distinct
residues and hence difference valuation $0$. Pulling back along their
common $m$ branches gives \eqref{epd:eq:first-mismatch}. In particular the
cosets in \eqref{epd:eq:fiber-main} are disjoint. The last assertion follows
from \cref{epd:lem:ideal}.
\end{proof}

\begin{corollary}[Exact finite cylinders]\label{epd:cor:cylinders}
For $x,y\in\BI(F)$ and $N\geq0$, the first $N$ itinerary symbols agree
if and only if $v(x-y)\geq N\kappa$.
\end{corollary}
\begin{proof}
A first mismatch $m<N$ gives valuation $m\kappa<N\kappa$. If the first
$N$ symbols agree, either the first mismatch is later, giving the
required inequality, or the itineraries are equal and their difference
belongs to $\Ik$.
\end{proof}

\begin{corollary}[The action on the fibers]\label{epd:cor:fiber-bijection}
$F$ maps $x_s+\Ik$ bijectively onto $x_{\sh s}+\Ik$. Every point of
$\BI(F)$ has exactly $d$ preimages in $\BI(F)$, one in each allowable
preceding-symbol fiber.
\end{corollary}
\begin{proof}
Forward inclusion is \eqref{epd:eq:fiber-growth}. Given a point in the
target fiber, apply $\Phi_{s_0}$. Its itinerary is $s$, by uniqueness
of the inverse branch, and it is the only preimage in that residue
class. There are $d$ residue classes and $F$ has degree $d$.
\end{proof}

This theorem identifies the information content of a full infinite
itinerary. It determines every finite $q$-scale but, in general, not a
point. The unresolved ambiguity is exactly an affine ideal, not an
unspecified smaller disk and not a choice of a formal completion.

\section{Periodic and preperiodic rigidity}
\label{epd:sec:periodic}
\begin{theorem}[All preperiodicity equations split]\label{epd:thm:periodic}
A point of $\BI(F)$ is preperiodic if and only if it is $x_s$ for an
eventually periodic word $s$. For every $m\geq0$ and $p\geq1$, the
polynomial
\begin{equation}\label{epd:eq:preperiodic-polynomial}
 F^{m+p}(X)-F^m(X)
\end{equation}
has exactly $d^{m+p}$ distinct roots, all in $K$ and all in $\BI(F)$.
They are precisely the centers with $\sh^{m+p}s=\sh^m s$. In particular
\[
 \#\Fix(F^p)=d^p.
\]
There are no preperiodic points outside $\OO$.
\end{theorem}
\begin{proof}
First consider a periodic word $s$ of period dividing $p$. Its center
is fixed by $F^p$. If $x=x_s+h$ with $h\ne0$ in $\Ik$ were also fixed,
then \eqref{epd:eq:fiber-growth} would imply simultaneously
\[
 v(F^p(x)-F^p(x_s))=v(h)-p\kappa
 \quad\hbox{and}\quad
 v(F^p(x)-F^p(x_s))=v(h),
\]
a contradiction. Hence the center is the unique periodic point in its
fiber. An eventually periodic point maps to a periodic point after
some $m$ iterations. Its difference from the corresponding center
never becomes zero, by \eqref{epd:eq:fiber-growth}, unless it was zero
initially. This proves the first assertion; the converse follows from
\eqref{epd:eq:shift-center}.

There are $d^{m+p}$ words satisfying $\sh^{m+p}s=\sh^m s$: prescribe
the first $m+p$ symbols, and repeat the last $p$ thereafter. They give
distinct roots of \eqref{epd:eq:preperiodic-polynomial}. That polynomial
has degree $d^{m+p}$, since $d\geq2$ and $p\geq1$. Thus they are all
its roots, and it splits with simple roots in $K$.

For completeness, if $v(x)<0$, the leading term of $P(x)$ strictly
dominates every lower-degree term. Consequently
\begin{equation}\label{epd:eq:escape-step}
 v(F(x))=d\,v(x)-\kappa<v(x)<0.
\end{equation}
Such a point never returns to $\OO$ and cannot be preperiodic, because
its successive valuations strictly decrease.
\end{proof}

\begin{proposition}[Exact multipliers]\label{epd:prop:multiplier}
For $x\in\BI(F)$ and $n\geq1$,
\[
 v\bigl((F^n)'(x)\bigr)=-n\kappa.
\]
At a root of \eqref{epd:eq:preperiodic-polynomial}, its derivative has
valuation $-(m+p)\kappa$.
\end{proposition}
\begin{proof}
At every integral orbit point, $P'$ has residue $P_0'(c_{s_j})\ne0$.
Thus $v(F')=-\kappa$, and the chain rule gives the first formula.
At a root $x$ of \eqref{epd:eq:preperiodic-polynomial}, factor its derivative
as
\[
 (F^m)'(x)\left((F^p)'(F^m(x))-1\right).
\]
The second factor has valuation $-p\kappa$, since that is negative and
$v(1)=0$. This proves the claimed valuation and directly confirms
simplicity.
\end{proof}

\begin{corollary}[Periodic counts and zeta series]\label{epd:cor:zeta}
The number of points of exact period $p$ is
\[
 N_p=\sum_{e\mid p}\mu(p/e)d^e,
\]
and the number of cycles of exact period $p$ is $N_p/p$. The formal
Artin--Mazur zeta series is
\[
 \exp\!\left(\sum_{p\geq1}\frac{\#\Fix(F^p)}p z^p\right)
 =\frac1{1-dz}.
\]
\end{corollary}
\begin{proof}
Every point fixed by $F^p$ has a unique exact period dividing $p$.
M\"obius inversion of $d^p=\sum_{e\mid p}N_e$ gives the first formula.
Each cycle has $p$ points. The last identity is the formal identity
$\sum_{p\geq1}(dz)^p/p=-\log(1-dz)$.
\end{proof}

The formulas themselves are classical full-shift formulas. Their role
here is to show what periodic data \emph{fails} to detect: the same
counts and multipliers persist whether the nonescaping set consists
of individual centers or of open infinitesimal fibers.

\section{The native topology and the scale quotient}
\label{epd:sec:topology}
Give $\Sig$ the product topology of the finite discrete alphabet.
Its cylinder of depth $N$ fixes the first $N$ symbols.

\begin{theorem}[Order-unit dichotomy]\label{epd:thm:topology}
The following conclusions hold.
\begin{enumerate}[label=\textup{(\alph*)},leftmargin=2em]
\item If $\kappa$ is an order unit, then $\Ik=0$ and
$s\mapsto x_s$ is a homeomorphism $\Sig\to\BI(F)$, conjugating the
one-sided shift to $F$. The set $\BI(F)$ is compact and perfect.
Periodic points are dense in it. The family $(F^n)_{n\geq0}$ is not
equicontinuous at any point of $\BI(F)$, even relative to that set.
\item If $\kappa$ is not an order unit, then $\BI(F)$ is clopen and noncompact in
$K$, every fiber $x_s+\Ik$ is open, and every subset of
$\CS:=\{x_s:s\in\Sig\}$ is closed and discrete in $K$. The itinerary
map $\pi:\BI(F)\to\Sig$ is continuous but is not a quotient map.
The native topological quotient by the fibers is discrete.
\item In the non-order-unit case the family $(F^n)_{n\geq0}$ is uniformly
equicontinuous on $\BI(F)$ in the native valuation uniformity.
\end{enumerate}
\end{theorem}
\begin{proof}
When $\kappa$ is an order unit, \cref{epd:cor:cylinders} and the order-unit property identify
the cylinder basis with a neighborhood basis on $\BI(F)$. The center
map and its inverse are continuous. Compactness and perfectness follow
from the elementary corresponding properties of the finite-alphabet
product. A periodic word extending any finite prefix proves density.
For failure of equicontinuity, fix any $x_s$ and any input neighborhood.
Choose a word $r$ sharing a sufficiently long prefix with $s$ but first
differing at a later index $m$. The input difference has valuation
$m\kappa$, arbitrarily large because $\kappa$ is an order unit; the difference after $m$
iterations has valuation $0$. The fixed output neighborhood $\{v>0\}$
cannot be maintained.

Now choose $\eta>n\kappa$ for every $n$. The fibers are open by
\cref{epd:lem:ideal}, so their union is open. Also
$\BI(F)=\bigcap_{n\geq0}(F^n)^{-1}(\OO)$ is closed, since $\OO$ is
closed and polynomials are continuous. Distinct centers have difference
valuation $m\kappa<\eta$. Thus a ball of threshold $\eta$ contains
at most one point of any subset $A\subseteq\CS$, proving discreteness.
To prove closedness, a point in the closure of $A$ has such a ball
meeting $A$ in exactly one point $a$. Unless the original point equals
$a$, a smaller ball excludes $a$, contradicting closure. Hence $A$ is
closed. In particular $\CS$ is an infinite closed discrete subset of
$\BI(F)$. If $\BI(F)$ were compact, so would be $\CS$, contradicting
the open cover of $\CS$ by its singletons. This proves noncompactness.

By \cref{epd:cor:cylinders}, every cylinder preimage is open, so $\pi$ is
continuous. Each individual fiber is open, whereas a singleton in
$\Sig$ is not open. Consequently $\pi$ is not a quotient map. The
quotient topology by open fibers is discrete.

Finally, fix an arbitrary desired output threshold $\gamma\in\Gamma$
and set
\[
 \delta=\max(\gamma,0)+\eta.
\]
If $x,y\in\BI(F)$ have $v(x-y)>\delta$, their difference lies in
$\Ik$, so they have the same itinerary by \cref{epd:thm:fibers}. Repeated
use of the exact forward metric yields
\[
 v(F^n(x)-F^n(y))=v(x-y)-n\kappa
 >\max(\gamma,0)+\eta-n\kappa>\gamma
\]
for every $n$. The same input threshold works for every pair and every
iterate, which is uniform equicontinuity.
\end{proof}

\begin{remark}[Why this is not a contradiction]\label{epd:rem:expansion}
The valuation of a nonzero difference within a fiber strictly decreases
by $\kappa$ at every step. Uniform equicontinuity nevertheless holds
because an input error can be placed beyond \emph{all} the scales that
this countable iteration can spend. A negative multiplier valuation
therefore does not by itself establish topological repulsion. This is
why the article does not identify its expanding points with classical
Julia points without naming a topology and a valuation scale.
\end{remark}

\subsection{The correct full-shift quotient}
\begin{definition}\label{epd:def:scaletop}
On $\BI(F)$, the \emph{$\kappa$-scale uniformity} is generated by
\[
 E_N=\{(x,y):v(x-y)\geq N\kappa\},\qquad N\in\N.
\]
Its indistinguishability relation is $x-y\in\Ik$.
\end{definition}
Each $E_N$ is an equivalence relation on $\BI(F)$, by the ultrametric
inequality. The family defines a uniformity even when it is not
Hausdorff.

\begin{proposition}\label{epd:prop:scalequotient}
The Hausdorff quotient of the $\kappa$-scale uniformity is uniformly
isomorphic to $\Sig$ with its cylinder uniformity, and the induced map
is the full shift. In particular its topological entropy is $\log d$.
In the non-order-unit case this is not the native topological quotient.
\end{proposition}
\begin{proof}
By \cref{epd:cor:cylinders}, $E_N$ is exactly equality of the first $N$
symbols. By \cref{epd:thm:fibers}, the common equivalence relation of all
$E_N$ is equality of the entire word, or equivalently difference in
$\Ik$. The quotient therefore has precisely the cylinder uniformity.
The shift identity follows from the definition of itinerary.

For entropy, there are exactly $d^N$ orbit names of length $N$ at the
resolution which distinguishes individual symbols. Refining to a
fixed cylinder length $r$ changes this count to $d^{N+r-1}$, without
changing the exponential rate. Hence the rate is $\log d$.
\end{proof}

\subsection{Compact invariant sets disappear at the finer scale}
\begin{theorem}[Compact-subsystem rigidity]\label{epd:thm:compact}
Suppose $\kappa$ is not an order unit. Every compact subset
$A\subseteq\BI(F)$ satisfying $F(A)\subseteq A$ is finite. Every point
of $A$ is a preperiodic center. Thus every compact forward-invariant
subsystem has topological entropy zero, although the scale quotient
in \cref{epd:prop:scalequotient} has entropy $\log d$.
\end{theorem}
\begin{proof}
The open fiber partition covers $A$. Compactness implies that only
finitely many fibers meet $A$. Since $F(A)\subseteq A$, the shift
orbit of each itinerary occurring in $A$ lies in a finite set. Each
such itinerary is therefore eventually periodic.

Suppose $x=x_s+h\in A$ with $h\ne0$. Choose $m\geq0$ and $p\geq1$
so that $\sh^{m+p}s=\sh^m s$, and put $c=x_{\sh^m s}$. For $j\geq0$
let $y_j=F^{m+jp}(x)\in A$. The corresponding center is always $c$,
and \cref{epd:thm:fibers} gives
\[
 v(y_j-c)=v(h)-(m+jp)\kappa.
\]
These valuations strictly decrease. For $j>i$, the two valuations
are unequal, so
\[
 v(y_j-y_i)=v(h)-(m+jp)\kappa\leq v(h)-m\kappa.
\]
The infinite set $\{y_j:j\geq0\}$ is uniformly separated. The same
ball argument used in \cref{epd:thm:topology} makes it closed and discrete
in $K$. It would be a closed subset of compact $A$, hence compact,
but an infinite discrete space is not compact. This contradiction
proves $h=0$ for every $x\in A$.

Thus $A\subseteq\CS$. Every subset of $\CS$ is discrete, so compactness
again forces $A$ to be finite. Forward invariance of a finite set makes
all its points preperiodic. A self-map of a finite space has entropy
zero.
\end{proof}

One must allow finite preperiodic tails in this theorem: the hypothesis
is $F(A)\subseteq A$, not $F(A)=A$. No conclusion that every point is
periodic is intended. There is also no conflict with entropy monotonicity
for compact factors: the native source space here is noncompact, and
the full shift is the quotient of a specifically coarsened uniformity.

\section{Scalar extension: new fibers, no new preperiodic points}
\label{epd:sec:extension}
Let $\Gamma\hookrightarrow\Delta$ be an order-preserving group
embedding, and use the induced coefficient-preserving Hahn embedding
\[
 K=k((t^\Gamma))\hookrightarrow L=k((t^\Delta)).
\]
The parameters $P,q$ are unchanged. The distinction between the two
ambient fields is real even if the polynomial coefficients occupy a
very small common subfield.

\begin{theorem}[Scalar-extension theorem]\label{epd:thm:extension}
The universal centers over $K$ and over $L$ are the same elements of
$K\subseteq L$. Moreover,
\begin{align}
 \BI(F;L)&=\coprod_{s\in\Sig}\bigl(x_s+\Ik(L)\bigr),\label{epd:eq:extension-fibers}\\
 \BI(F;L)\cap K&=\BI(F;K),\qquad
 \Ik(L)\cap K=\Ik(K).\label{epd:eq:intersection-extension}
\end{align}
Every periodic and preperiodic point over \emph{any field extension}
of $K$ already belongs to $K$. If $\kappa$ is an order unit of $\Gamma$ but
not of $\Delta$, the following additional conclusions hold:
\begin{enumerate}[label=\textup{(\roman*)},leftmargin=2em]
\item $\BI(F;K)$ is the compact center set in the intrinsic topology of
$K$, but it is closed and discrete in the intrinsic topology of $L$.
\item Every point of $\BI(F;L)\setminus\BI(F;K)$ is transcendental over
$K$.
\item The restriction of the topology of $L$ to $K$ is not the intrinsic
valuation topology of $K$.
\end{enumerate}
\end{theorem}
\begin{proof}
The two evaluations of each universal formal series agree
coefficientwise. All supports and all finite coefficient sums are
preserved by the ordered exponent embedding. Apply
\cref{epd:thm:fibers} in $L$ to get \eqref{epd:eq:extension-fibers}.
The value comparison with $n\kappa$ is unchanged on $K$, proving the
ideal intersection. Integral membership and every polynomial iterate
are also unchanged on $K$, which proves the nonescaping intersection.

For any field extension $M/K$, a preperiodic element is a root of
$F^{m+p}-F^m$ for some $m,p$. By \cref{epd:thm:periodic}, that polynomial
already splits into distinct linear factors in $K[X]$. A root in $M$
therefore equals one of those roots in $K$. This part does not require
$M$ to be a Hahn field or even to carry a valuation.

Assume now that the order-unit property is lost. The first conclusion is
\cref{epd:thm:topology} in the two fields. For the second, write a new point
as $x_s+h$ with $h\in\Ik(L)\setminus\{0\}$. Because $\kappa$ is
an order unit of $\Gamma$, $v(h)$ exceeds every element of $\Gamma$.
Suppose a nonzero polynomial $\sum_{j=0}^r b_jh^j=0$, with $b_j\in K$,
existed. Let $j_0$ be its least index with $b_{j_0}\ne0$. For $j>j_0$
with $b_j\ne0$,
\[
 v(b_jh^j)-v(b_{j_0}h^{j_0})
 =v(b_j)-v(b_{j_0})+(j-j_0)v(h)>0.
\]
Thus the $j_0$ term is the unique term of least valuation and cannot
cancel. This contradiction proves $h$, hence $x_s+h$, transcendental
over $K$.

Finally, on the center set the topology from $K$ is a non-discrete
Cantor topology, whereas that from $L$ is discrete. They cannot agree
on all of $K$.
\end{proof}

\begin{example}[A minimal extension that changes the dynamics]\label{epd:ex:minimal-extension}
Take
\[
 K=\C((q)),\qquad
 L=\C((t^{\Q\oplus_{\rm lex}\Q})),
 \qquad q\longmapsto t^{(0,1)},\qquad F=q^{-1}(X^2-1).
\]
Here $K$ is identified with the Hahn subfield with exponent subgroup
$\{0\}\oplus\Z$. The integral nonescaping set in $K$ is the binary
center set. In $L$, each center acquires the fiber $x_s+I_{(0,1)}$;
for example $x_s+t^{(1,0)}$ is a new nonescaping point. None of these
new points is algebraic over $K$. Nevertheless every solution of every
preperiodicity equation was already in $K$.
\end{example}

The theorem separates two operations that should not be conflated:
extending the field and taking a symbolic or topological model of the
nonescaping set. The algebraic equations for all finite preperiodic
orbits are stable under extension, while the infinite itinerary
realization problem is not pointwise stable.

\section{Support-controlled deformations and invisible changes}
\label{epd:sec:deformation}
Keep $q$ fixed. Let $\widetilde P=P+D$, where
\[
 D(X)=\sum_{j=0}^d d_jX^j,\qquad d_j\in\mm,
 \qquad \widetilde F=q^{-1}\widetilde P.
\]
Both polynomials have the same labeled simple reduction $P_0$.
Write $x_s$ and $\widetilde x_s$ for their respective universal
centers.

\begin{theorem}[Uniform deformation estimate]\label{epd:thm:deformation}
If $D\ne0$ and $\theta=\min_{d_j\ne0}v(d_j)>0$, then for every
itinerary $s$,
\begin{equation}\label{epd:eq:deformation-bound}
 v(\widetilde x_s-x_s)\geq\theta.
\end{equation}
The supports of all differences belong to the monoid generated by
\[
 \supp(q)\cup\bigcup_j\supp(e_j)\cup\bigcup_j\supp(d_j).
\]
Let $D_\theta(X)=\sum_j(\coeff{\theta}{d_j})X^j\in k[X]$. The first
possible correction is exactly
\begin{equation}\label{epd:eq:leading-response}
 \coeff{\theta}{(\widetilde x_s-x_s)}
 =-\frac{D_\theta(c_{s_0})}{P_0'(c_{s_0})}.
\end{equation}
If the numerator is nonzero, equality holds in
\eqref{epd:eq:deformation-bound}. If it vanishes, cancellation is allowed;
no equality is asserted.
\end{theorem}
\begin{proof}
Introduce further formal variables $D_0,\ldots,D_d$. In the
finite-variable power-series ring in $Q,E,D$, the difference
\[
 X_s(Q,E+D)-X_s(Q,E)
\]
belongs to $(D_0,\ldots,D_d)$. Indeed its specialization at $D=0$ is
zero, so each monomial has at least one $D$ factor. Apply
\cref{epd:lem:evaluation} at $Q=q$, $E=e$, $D=d$. The two evaluations
are exactly $\widetilde x_s$ and $x_s$, by compatibility with formal
substitution. Each difference is a finite sum $\sum_jd_jB_{s,j}$ with
$B_{s,j}\in\OO$. This proves the uniform valuation bound. The support
statement is the same evaluation lemma.

Put $\Delta_s=\widetilde x_s-x_s$. Subtract the two shift equations:
\begin{equation}\label{epd:eq:variation-eq}
 P(x_s+\Delta_s)-P(x_s)+D(x_s+\Delta_s)
       =q\Delta_{\sh s}.
\end{equation}
The right side has valuation at least $\kappa+\theta>\theta$.
Every term of degree at least two in $\Delta_s$ has valuation greater
than $\theta$. At exponent $\theta$, the first difference on the left
is therefore $P_0'(c_{s_0})\coeff{\theta}{\Delta_s}$. At the same
exponent, the second term is $D_\theta(c_{s_0})$. Extracting this
coefficient in \eqref{epd:eq:variation-eq} proves
\eqref{epd:eq:leading-response}.
\end{proof}

\begin{remark}[Why a formal ideal argument is necessary]\label{epd:rem:formal-ideal}
The threshold $\theta$ may exceed all $n\kappa$. Then estimates for
finite nested approximants with error valuation $n\kappa$ do not
justify a claim about the $\theta$ coefficient of an alleged limit.
The proof instead establishes divisibility by the perturbation ideal
\emph{before} evaluating. This is the same support-versus-topology
issue that required the universal construction of the centers.
\end{remark}

\begin{corollary}[Deformations invisible at all dynamical orders]
\label{epd:cor:flatdeform}
If every coefficient $d_j$ belongs to $\Ik$, then
\[
 \BI(F)=\BI(\widetilde F),\qquad
 \pi_F^{-1}(s)=\pi_{\widetilde F}^{-1}(s)\quad(s\in\Sig).
\]
The two maps induce the same full shift on the scale quotient. For
any nonzero constant perturbation $D=\varepsilon\in\Ik$, however,
$F$ and $\widetilde F$ have no common fixed point.
\end{corollary}
\begin{proof}
There are finitely many coefficients. If $D\ne0$, their minimum
nonzero valuation still exceeds every $n\kappa$.
\Cref{epd:thm:deformation} therefore gives
$\widetilde x_s-x_s\in\Ik$ for every $s$. The fiber formula makes
the labeled cosets identical. Both quotient actions shift the labels.
Alternatively, $D(x)/q\in\Ik$ for integral $x$ because $q\Ik=\Ik$.
If a common fixed point existed for a nonzero constant $D$, it would
satisfy both $P(x)=qx$ and $P(x)+\varepsilon=qx$, which is impossible.
\end{proof}

Thus even the entire infinite symbolic partition, together with all
finite-scale data, does not determine the actual fixed points in a
workspace in which $\kappa$ is not an order unit. The correction formula shows exactly where the
missing information first appears.

\section{Integral nonescape versus valuation-bounded orbits}
\label{epd:sec:bounded}
Define the other natural orbit locus by
\[
 \BV(F)=\{x\in K:\text{there exists }\beta\in\Gamma
                \text{ with }v(F^n(x))\geq\beta\text{ for all }n\geq0\}.
\]
Always $\BI(F)\subseteq\BV(F)$. The opposite inclusion is often
silently assumed in a rank-one argument. Here it has an exact
order-unit criterion.

For $x\ne0$, set $\alpha(x)=\max(0,-v(x))$, and put $\alpha(0)=0$.
Thus $\alpha$ measures only how far the initial point lies outside
the integral disk.

\begin{theorem}[Exact bounded-orbit locus]\label{epd:thm:bounded}
For $F=q^{-1}P$ as in \eqref{epd:eq:polynomial},
\begin{equation}\label{epd:eq:bounded-formula}
 \BV(F)=\BI(F)\ \cup\
 \{x\in K:\kappa+\alpha(x)\text{ is not an order unit in }\Gamma\}.
\end{equation}
Consequently:
\begin{enumerate}[label=\textup{(\alph*)},leftmargin=2em]
\item If $\kappa$ is an order unit, $\BV(F)=\BI(F)$.
\item If $\kappa$ is not an order unit, then
$\BV(F)=\{x:\kappa+\alpha(x)\text{ is not an order unit}\}$.
In this case the bounded-orbit locus is independent of $P$ within the
specified family, and it contains all of $\OO$.
\item If $\Gamma$ has no order unit, then $\BV(F)=K$.
\end{enumerate}
\end{theorem}
\begin{proof}
If an orbit ever leaves $\OO$, write its first negative valuation as
$-a$, with $a>0$. If the initial point was integral, then $0<a\leq\kappa$:
the preceding point was integral, $P$ has integral coefficients, and
$v(F(z))\geq-\kappa$ for integral $z$. If the initial point was not
integral, we may start at time zero with $a=\alpha(x)$.

Once negative, valuations obey \eqref{epd:eq:escape-step}. After a further
$n$ steps they are
\begin{equation}\label{epd:eq:outside-recurrence}
 -d^n a-(1+d+\cdots+d^{n-1})\kappa.
\end{equation}
For $n\geq1$, the positive quantity whose negative occurs here lies
between $d^{n-1}(a+\kappa)$ and $d^n(a+\kappa)$. Its values are
cofinal in $\Gamma_{>0}$ exactly when $a+\kappa$ is an order unit.
Indeed, if that element is an order unit, powers of $d$ dominate any finite
multiple of it; otherwise a single $\eta$ exceeding all its integer
multiples bounds the whole sequence. A finite initial orbit segment
does not affect boundedness.

When the initial point is integral and later exits, $a+\kappa$ lies
between $\kappa$ and $2\kappa$, so it is an order unit exactly when $\kappa$
is. Otherwise $a+\kappa=\alpha(x)+\kappa$. Points which never leave
are already in $\BI(F)$. This proves \eqref{epd:eq:bounded-formula}.
The three consequences follow directly; for (b), all points of
$\BI(F)$ have $\alpha=0$, so the first set in the union is redundant.
\end{proof}

\subsection{A coarsened valuation-ring interpretation}
\label{epd:subsec:coarsened}
Suppose $\Gamma$ has an order unit $u$, and define
\[
 H=\{\gamma\in\Gamma:n|\gamma|<u\text{ for every integer }n\geq1\}.
\]
This is the maximal proper convex subgroup of $\Gamma$. In particular,
a positive element belongs to $H$ exactly when it is not an order unit.
For clarity, the relevant elementary argument is included. Closure
under addition follows by applying the condition to $2n|\gamma|$ and
$2n|\delta|$: their sum is less than $2u$, hence
$n(|\gamma|+|\delta|)<u$. Convexity is immediate. If $\gamma>0$ is not
an order unit, some $\eta$ exceeds all $n\gamma$, while $\eta\leq Nu$ for
some $N$. Applying the bound to $Nn\gamma$ gives $n\gamma<u$.
Conversely that last condition prevents $\gamma$ from being an order unit. Any convex subgroup
containing an order unit is all of $\Gamma$, proving maximality.
(Placement note: $H$ is the kernel of the order-preserving homomorphism
$\Gamma\to\R$ of \texttt{ent:lem:coarsening} in
\texttt{entire-functions-\allowbreak at-arbitrary-rank}, with $h=u$. That report
states the lemma under its standing divisibility hypothesis. The argument
just given does not use divisibility.)

\begin{corollary}\label{epd:cor:coarsebounded}
If $\Gamma$ has an order unit but $\kappa$ is not one, then
\[
 \BV(F)=\{x\in K:v(x)+H\geq0\text{ in }\Gamma/H\},
\]
the valuation ring of the coarsened valuation. Relative to this
coarsening, $F$ has integral coefficients and unit leading coefficient.
\end{corollary}
\begin{proof}
Here $\kappa\in H$, and for $\alpha\geq0$ the sum $\kappa+\alpha$
is not an order unit exactly when $\alpha\in H$. The formula follows from
\cref{epd:thm:bounded}. Also $v(q)\in H$, so $q$ is a unit for the
coarsened valuation. The coefficients of $P$ are integral and its
leading coefficient is a unit already for the original valuation.
\end{proof}

\begin{example}[Escaping the disk without an unbounded orbit]\label{epd:ex:escape}
With $\Gamma=\Q\oplus_{\mathrm{lex}}\Q$ and $\kappa=(0,1)$,
$H=\{0\}\oplus\Q$. Therefore $\BV(F)$ consists of those elements
whose valuation has nonnegative first coordinate. For
$x=t^{(0,-1)}$ the integral disk has already been left, and
\eqref{epd:eq:outside-recurrence} gives increasingly negative second
coordinates. All valuations remain greater than $(-1,0)$, so the
orbit is still valuation-bounded. This is not a failure of the
recurrence; it is the distinction between a direction that is not an order unit and
a genuinely unbounded sequence in the full ordered group.
\end{example}

\section{A general polynomial theorem when there is no order unit}
\label{epd:sec:general}
The absence of an order unit has a consequence beyond the
simple-reduction family. The following quantitative estimate explains
why the native fine topology can conceal every finite-degree
polynomial expansion.

\begin{proposition}[An orbit and difference budget]\label{epd:prop:budget}
Let $f(X)=\sum_{j=0}^d a_jX^j\in K[X]$ have degree $d\geq1$.
Choose $A>0$ with $v(a_j)\geq-A$ for all nonzero coefficients, and put
$D_A=\{z:v(z)\geq-A\}$. For $x,y\in D_A$ and $n\geq0$,
\begin{align}
 v(f^n(x))&\geq-(d+1)^n A,\label{epd:eq:orbitbudget}\\
 v(f^n(x)-f^n(y))&\geq v(x-y)-\bigl((d+1)^n-1\bigr)A.
 \label{epd:eq:differencebudget}
\end{align}
\end{proposition}
\begin{proof}
Set $M_n=(d+1)^nA$. If $v(z)\geq-M_n$, every polynomial term has
valuation at least $-A-dM_n\geq-(d+1)M_n=-M_{n+1}$. This proves
\eqref{epd:eq:orbitbudget} by induction.

For $z,w$ with valuations at least $-M_n$, the divided difference
\[
 \frac{f(z)-f(w)}{z-w}
 =\sum_{j=1}^d a_j\sum_{\ell=0}^{j-1}z^\ell w^{j-1-\ell}
\]
(where the right side also defines it when $z=w$) has valuation at
least $-A-(d-1)M_n\geq-dM_n$. Apply this with
$z=f^n(x)$ and $w=f^n(y)$. Accumulating the possible losses gives
\[
 d\sum_{j=0}^{n-1}M_j=((d+1)^n-1)A,
\]
which is \eqref{epd:eq:differencebudget}. Cancellation can only increase
these valuations.
\end{proof}

\begin{theorem}[Universal native equicontinuity criterion]
\label{epd:thm:general}
The following statements are equivalent for $K=k((t^\Gamma))$:
\begin{enumerate}[label=\textup{(\roman*)},leftmargin=2em]
\item $\Gamma$ has no order unit.
\item Every orbit of every polynomial in $K[X]$ is valuation-bounded.
\item For every polynomial $f$, its family of iterates is locally
equicontinuous at every point of $K$, in the native valuation uniformity.
\end{enumerate}
Under (i), the family is in fact uniformly equicontinuous on every
$D_A$ for which the coefficient bounds in \cref{epd:prop:budget} hold.
\end{theorem}
\begin{proof}
Assume (i). For the $A$ in \cref{epd:prop:budget}, choose
$\eta>rA$ for every ordinary integer $r$. Then
\eqref{epd:eq:orbitbudget} bounds every orbit starting in $D_A$ below by
$-\eta$. Every individual starting point lies in some such $D_A$,
proving (ii).

For any requested output threshold $\gamma$, choose
$\delta=\max(\gamma,0)+\eta$. If $x,y\in D_A$ and
$v(x-y)>\delta$, then \eqref{epd:eq:differencebudget} gives
$v(f^n(x)-f^n(y))>\gamma$ for every $n$. This proves uniform
equicontinuity on $D_A$. Such a disk is a neighborhood of each of its
points: the ball $\{y:v(y-x)>-A\}$ is contained in it. Thus (iii)
follows. Constant polynomials need no estimate and satisfy the same
conclusions.

Conversely, suppose $u>0$ is an order unit. The linear polynomial
$f(X)=t^{-u}X$ has the orbit $f^n(1)=t^{-nu}$, whose valuations have
no lower bound. Hence (ii) fails. Its iterates are not locally
equicontinuous at zero: given any input threshold $\delta$, choose
$N$ with $Nu>\delta$ and take $h=t^{Nu}$. Then $h$ belongs to that
input neighborhood, but $v(f^{N+1}(h))=-u<0$. Even the fixed output
threshold $0$ cannot be maintained. Thus (iii) fails as well.
\end{proof}

\begin{definition}[A deliberately named affine notion]\label{epd:def:affine-julia}
For a polynomial $f$, define $J_{\mathrm{aff},v}(f)$ to be the set of
points of $K$ having no native open neighborhood on which all iterates
of $f$ are equicontinuous as maps to the additive valued field $K$.
\end{definition}

\begin{corollary}\label{epd:cor:empty-affine}
If $\Gamma$ has no order unit, then $J_{\mathrm{aff},v}(f)=\varnothing$
for every polynomial $f$.
\end{corollary}
\begin{proof}
This is exactly statement (iii) of \cref{epd:thm:general}.
\end{proof}

The subscript is essential. Classical non-Archimedean Fatou and Julia
sets are ordinarily defined on a projective line using a spherical
metric, or on a Berkovich space, not by the specifically affine
replacement; compare Definitions~4.15--4.16 of \cite{Benedetto}.
We claim neither that every Berkovich Julia set is empty nor that a
rank-one Julia set disappears in its own topology. The corollary is a
precise obstruction to one naive transplantation of the definition.

\begin{example}[No uncountable-cofinality hypothesis is needed]\label{epd:ex:no-order-unit}
Let $\Gamma=\bigoplus_{j\geq0}\Q e_j$, ordered by the largest index
having a nonzero coefficient, with its sign deciding the order.
Then $e_{j+1}>ne_j$ for every positive integer $n$, and every group
element has finite support. The group is countable and has no order
unit. All conclusions of \cref{epd:thm:general} hold. This is a finite
polynomial growth-budget phenomenon, not a consequence of assuming
uncountable cofinality for the entire value group.
(Placement note: here, as in \texttt{ent:rem:one-name}, cofinality means
$\operatorname{cf}(\Gamma)$, which is $\aleph_0$ for this group. As an
ordered group it is the group $\Gamma_\infty=\bigoplus_{j\geq0}\Q\,\omega^j$
of \texttt{ent:subsec:gamma-infinity}, via $e_j=\omega^j$.)
\end{example}

\section{The surreal and surcomplex specialization}
\label{epd:sec:surreal}
\subsection{What is imported from surreal theory}
Conway normal form identifies surreal numbers with set-supported
well-based sums of real coefficients times monomials $\omega^a$.
The monomial law $\omega^{a+b}=\omega^a\omega^b$ identifies a
set-sized ordered additive subgroup $\Gamma\subseteq\No$ with a
Hahn monomial group by $t^\gamma=\omega^{-\gamma}$. Thus
$\R((t^\Gamma))$ embeds into $\No$, and coefficientwise
complexification embeds $\C((t^\Gamma))$ into $\SC$.
This is classical normal-form theory, not a theorem first proved here;
see Gonshor \cite{Gonshor} and the explicit Hahn description in
Berarducci--Mantova, Sections~2.3, 2.4, and 2.6 \cite{BM}.

The valuation convention remains $v(\omega^{-\gamma})=\gamma$.
For a surcomplex number written in normal form, coefficients are in
$\C$ and supports are still sets. Every finite list of parameters is
contained in a common set-sized Hahn workspace obtained by taking the
additive group generated by its supports and $1$, so that the workspace
value group is nonzero even for purely ordinary parameters. All infinite sums used to
construct the centers have the common set-sized support certificate
\eqref{epd:eq:common-support}.

There is consequently no need to posit a set containing all surreals,
or to sum over a proper class of exponents. Statements about an
arbitrary additional point are proved after enlarging the workspace to
include that point's support. The same center is obtained after every
such enlargement by \cref{epd:thm:extension}.

\subsection{An explicit quadratic theorem inside the actual field}
Set
\begin{equation}\label{epd:eq:surreal-F}
 q=\omega^{-1},\qquad F(x)=\omega(x^2-1),\qquad
 I=\{h:v(h)>n\text{ for every }n\in\N\}.
\end{equation}
Here $h$ may be surreal or surcomplex. In the real case, $I$ consists
of numbers smaller in magnitude than every positive ordinary power
of $\omega^{-1}$; the valuation definition is the one used in the
proofs. It contains $\omega^{-\omega}$.

\begin{theorem}[A surcomplex full shift with infinitesimal fibers]
\label{epd:thm:surreal}
There is a canonical real surreal number
\[
 x_s=\sum_{n\geq0}a_n(s)\omega^{-n}
 \qquad\bigl(s\in\{+1,-1\}^{\N}\bigr)
\]
for every binary itinerary. The locus whose every iterate is a finite surcomplex number is
\begin{equation}\label{epd:eq:surreal-fibers}
 \{x\in\SC:F^n(x)\text{ is finite for every }n\}
   =\coprod_s(x_s+I).
\end{equation}
Its restriction to $\No$ is obtained by using real $h\in I$.
Every actual surcomplex preperiodic point of $F$ is one of the real
centers with eventually periodic itinerary. For $h=\omega^{-\omega}$,
\begin{equation}\label{epd:eq:omega-error}
 v\bigl(F^n(x_s+h)-x_{\sh^n s}\bigr)=\omega-n
 \qquad(n\in\N).
\end{equation}
In the full fine topology the fibers and their union are clopen.
Every subset of the center set is closed and discrete. The itinerary map
is continuous but not a quotient map to the product-topological full shift;
its native quotient is discrete. The iterates are uniformly equicontinuous
on the nonescaping locus. Every compact set-sized forward-invariant subset
is finite and consists of preperiodic centers.
\end{theorem}
\begin{proof}
Apply the constant-coefficient construction in the workspace with
exponent group $\Z$. This produces every $x_s$ and its strong series.
For any proposed point or perturbation, enlarge to a set-sized exponent
group containing its support and $1$. The exact-fiber theorem gives
the same formula there, and the center is unchanged. Finite surreal
and surcomplex elements are exactly the elements of nonnegative
natural valuation, so the left side of \eqref{epd:eq:surreal-fibers} is
precisely the integral nonescaping condition.

The class $I$ is nonzero because $\omega>n$ for all $n$. The fiber
formula and \eqref{epd:eq:omega-error} follow from
\cref{epd:thm:fibers}. All coefficients of the binary centers are real.
Any preperiodic surcomplex point is a root of a polynomial which already
splits into distinct real linear factors by \cref{epd:thm:periodic}, so it
is one of those real centers. For topology, the proofs in
\cref{epd:thm:topology,epd:thm:compact} only use balls, a fixed dominating
valuation, and set-indexed compactness; they apply verbatim.
\end{proof}

\begin{remark}[Fine topology versus a workspace topology]\label{epd:rem:fine-vs-workspace}
The natural valuation topology of the full surreal field is its order
(fine) topology: every positive radius exceeds some monomial radius,
and increasing a monomial valuation gives a radius strictly smaller
than any prescribed positive element. In the surcomplex field the
same statement holds with modulus balls, or equivalently with the
product of the two real-coordinate topologies. These assertions also
follow directly by comparing a leading monomial with its nonzero real
or complex coefficient.

The restriction of this full topology to $\R((\omega^{-1}))$ is not
that field's intrinsic $\omega^{-1}$-adic topology. In the latter, the
center set is a Cantor set. In the former, the same set is uniformly
discrete, since a radius with valuation $\omega$ separates every pair
of centers. These are two different topologies on the same algebraic
objects, not two incompatible conclusions.
\end{remark}

\begin{remark}[Placement note: the set-level clauses are special cases]
\label{epd:rem:discrete-special}
The manuscript did not cite \emph{Surcomplex Analysis}
(\texttt{docs/surcomplex/analysis}). That report's fine topology
(\texttt{a:def:fine}) is the full topology used in
\cref{epd:thm:surreal,epd:rem:fine-vs-workspace}. Its \texttt{a:thm:discrete}
proves that every \emph{set} $A\subseteq\No[\mathrm i]$ is closed and
discrete in that topology, and that one radius separates all pairs of
distinct points of~$A$. The following consequences use nothing from this
report except, in the second item, the identification of the preperiodic
points.
\begin{itemize}[leftmargin=1.4em]
\item Every subset of the center set $\{x_s\}$ is closed and discrete,
because that set is indexed by the set $\{+1,-1\}^{\N}$.
\item Every compact set-sized subset is finite. In particular, every compact
forward-invariant set-sized subset of the locus \eqref{epd:eq:surreal-fibers}
is finite. Its points have finite orbits, so they are preperiodic, and by
\cref{epd:thm:surreal} they are real centers.
\item The uniform separating radius of \texttt{a:thm:discrete}(ii) is
realized explicitly by the radius of valuation $\omega$ in
\cref{epd:rem:fine-vs-workspace}. The
comparison made in that remark is an instance of
\texttt{a:rem:notvaluation}: the fine topology induced on a set-sized Hahn
subfield is strictly finer than the subfield's own valuation topology.
\end{itemize}
The set-level discreteness and compactness clauses in the last sentence of
\cref{epd:thm:surreal} are therefore special cases of
\texttt{a:thm:discrete}. The substance of \cref{epd:thm:topology}(b) and
\cref{epd:thm:compact} lies in a set-sized workspace $k((t^\Gamma))$ with
its intrinsic valuation topology. That topology is not discrete, and in the
order-unit case the center set is a Cantor set (\cref{epd:thm:topology}(a)). Some
statements concern the fibers $x_s+I$, which are proper classes, or whole
fine balls. These are the openness of the fibers, the failure of $\pi$ to
be a quotient map, uniform equicontinuity on the locus
\eqref{epd:eq:surreal-fibers}, and \cref{epd:cor:surreal-all}. None of them
is a consequence of \texttt{a:thm:discrete}.
\end{remark}

\subsection{Exact coefficients and the first visible digits}
For the quadratic family the inverse branches are
\[
 \Phi_{\varepsilon}(w)
 =\varepsilon\sqrt{1+qw}
 =\varepsilon\sum_{r\geq0}\binom{1/2}{r}q^rw^r,
 \qquad \varepsilon\in\{+1,-1\}.
\]
The root is the formal root with constant coefficient $1$. Its Hahn
evaluation is justified by the positive valuation of $qw$.
Write $a_n(s)$ for the coefficient of $q^n$ in $x_s$. Comparing
coefficients in $x_s^2-1=qx_{\sh s}$ gives
\begin{equation}\label{epd:eq:binary-recursion}
 a_0(s)=s_0,\qquad
 a_n(s)=\frac{a_{n-1}(\sh s)-
                 \sum_{j=1}^{n-1}a_j(s)a_{n-j}(s)}{2s_0}
 \quad(n\geq1).
\end{equation}
In particular,
\begin{equation}\label{epd:eq:first-coeffs}
 x_s=s_0+\frac{s_0s_1}{2}q+
       \left(\frac{s_0s_1s_2}{4}-\frac{s_0}{8}\right)q^2+\cdots.
\end{equation}
All coefficients are rational. The coefficient at order $n$ sees
exactly a finite prefix, although the resulting Hahn point carries
an arbitrary infinite word.

For the constant positive itinerary, the fixed point is
\begin{align}\label{epd:eq:fixed-plus}
 x_+&=\frac q2+\sqrt{1+\frac{q^2}{4}}\\
 &=1+\frac q2+\frac{q^2}{8}-\frac{q^4}{128}
     +\frac{q^6}{1024}-\frac{5q^8}{32768}+\cdots.\nonumber
\end{align}
The other fixed point is $x_-=q/2-\sqrt{1+q^2/4}$.
The point $x_++\omega^{-\omega}$ is not periodic, but every one of its
iterates has residue $+1$. Formula \eqref{epd:eq:omega-error} records its
exact deviation from the fixed point at every finite time.

A deformation gives another explicit check. Replace $P=X^2-1$ by
$\widetilde P=X^2-1+\varepsilon$, with
$\varepsilon=\omega^{-\omega}$. Then
\[
 \widetilde x_+=\frac q2+
       \sqrt{1-\varepsilon+\frac{q^2}{4}},
 \qquad
 [\omega^{-\omega}](\widetilde x_+-x_+)=-\frac12.
\]
The bracket denotes the normal-form coefficient of that monomial.
The two maps have exactly the same labeled nonescaping fibers, but no
common fixed point. This is the simplest instance of
\cref{epd:cor:flatdeform}.

\subsection{Polynomial iteration in the full fine topology}
For every positive surreal $A$ there is a surreal $\eta$ exceeding all
ordinary integer multiples of $A$, for example $\eta=\omega A$.
The quantitative estimates in \cref{epd:prop:budget} therefore apply to
any fixed polynomial with surreal or surcomplex coefficients and any
set of inputs in the specified ball. All relevant parameters and
$\eta$ can first be placed in a common set-sized workspace.

\begin{corollary}\label{epd:cor:surreal-all}
Every orbit of every finite-degree polynomial over $\No$ or $\SC$ is
bounded by some surreal modulus bound. The family of iterates of each
such polynomial is locally equicontinuous everywhere in the full fine
uniformity. Its explicitly affine equicontinuity obstruction
$J_{\mathrm{aff},v}$ is empty.
\end{corollary}
\begin{proof}
For a fixed polynomial and an initial ball, choose the bound $A$ in
\cref{epd:prop:budget} and then a larger scale $\eta$ dominating all $nA$.
The proof of \cref{epd:thm:general} uses only this particular pair of
scales. For each pair of inputs, enlarge the set-sized workspace to
contain their supports as well. The same $A$, $\eta$, and input
threshold work in every such enlargement, so the estimate is uniform
on the entire fine ball, even though that ball is not itself set-sized.
A lower valuation bound gives a modulus bound
by taking a still larger monomial to absorb every possible finite
leading coefficient. The same comparison of monomial balls
and modulus balls transfers equicontinuity to the fine uniformity.
\end{proof}

For the explicit map $F(x)=\omega(x^2-1)$, the orbit of $0$ satisfies
\[
 v(F^n(0))=-(2^n-1)\qquad(n\geq1).
\]
It leaves the integral disk immediately, yet every iterate has modulus
less than $\omega^\omega$. Thus $0\in\BV(F)\setminus\BI(F)$ in the
full surreal field. More simply, the ordinary sequence $2^{2^n}$ is
bounded by an infinite surreal number although it is unbounded over
$\R$. The content needed
for polynomial equicontinuity is stronger than mere existence of a
bound for one sequence: \eqref{epd:eq:differencebudget} provides a single
input precision that works for every iterate and every pair of inputs
in a whole prescribed ball. The corollary still does not assert that
classical or Berkovich Julia sets are empty in their own topologies.

\section{Verification, reproducibility, and formalization}
\label{epd:sec:verification}
The accompanying program \texttt{code/verify.py} uses Python's standard
library only, with exact rational coefficients and integer exponent
pairs. It is a finite consistency test of the displayed algebra, not
a substitute for the proofs of infinite support or compactness.

The recorded run performed \textbf{38,139 exact checks}, all passing.
The principal counts are shown in \cref{epd:tab:checks}. Every count is
computed by the shipped program; the full record is in
\texttt{data/verification.json}.

\begin{table}[ht]
\centering\small
\caption{Exact finite checks in the included run.}\label{epd:tab:checks}
\begin{tabularx}{\linewidth}{@{}Xr@{}}
\toprule
Check & Number\\\midrule
Two independent center constructions, all binary words of length $8$ & 256\\
Itinerary equation through the available truncation order & 256\\
Displayed first three coefficients & 256\\
Every unordered pair of length-$8$ words: exact first mismatch order & 32,640\\
Periodic identities, all binary words of periods $1$ through $6$ & 126\\
M\"obius exact-period counts & 6\\
Two-parameter identity, ideal divisibility, and first response & 96\\
Displayed positive fixed-point expansion & 1\\
Rank-two error valuation and leading coefficient, times $0$ through $10$ & 22\\
Exterior valuation recurrence versus its closed formula & 4,480\\\midrule
Total & 38,139\\\bottomrule
\end{tabularx}
\end{table}

\subsection{The two coefficient algorithms}
The first algorithm composes the inverse branches
$\varepsilon\sqrt{1+qw}$ using a truncated binomial evaluator. The
second uses \eqref{epd:eq:binary-recursion} directly, constructing the tail
coefficients before the coefficients of the current word. These are
different algorithms, though both implement the same defining
polynomial equation; their agreement is not independent foundational
verification.

For a requested coefficient order $N$, only $N+1$ symbols are needed.
The direct recursive implementation has polynomial arithmetic cost
in $N$; with straightforward convolutions and all required suffixes,
$O(N^3)$ rational arithmetic operations suffice. Coefficient bit growth
is not included in this arithmetic-operation bound. Nothing in this
finite computation decides membership in $I_\kappa$, which quantifies
over all finite $q$-orders.

The bivariate checks work in $\Q[q,\varepsilon]$ truncated by total
degree. They verify $x_s^2-1+\varepsilon=qx_{\sh s}$, specialization
at $\varepsilon=0$, and the coefficient
$[q^0\varepsilon^1](\widetilde x_s-x_s)=-s_0/2$. The rank-two checks
track the formal difference recurrence
\[
 h_{n+1}=q^{-1}(2x_+h_n+h_n^2)
\]
in powers of an additional symbol $h$, with lexicographically ordered
exponents $(\deg_h,\deg_q)$. Their leading terms are
$2^n hq^{-n}$, of valuation $(1,-n)$. Discarded higher $h$ degrees
cannot affect the degree-$1$ leading term. These checks do not pretend
to turn an approximate fixed-point truncation into an exact Hahn
fixed point; the exact point was proved in \eqref{epd:eq:fixed-plus}.

\subsection{How to reproduce the package}
\label{epd:subsec:reproduce}
The manuscript's instructions, for its delivered layout, were: from the
package directory, run
\begin{lstlisting}
python code/verify.py --output data/verification-rerun.json
pdflatex -interaction=nonstopmode -halt-on-error article.tex
pdflatex -interaction=nonstopmode -halt-on-error article.tex
pdflatex -interaction=nonstopmode -halt-on-error article.tex
\end{lstlisting}
The source has an internal bibliography and needs no BibTeX run.
\texttt{make all} performs the verification and PDF build. Python
3.10 or later is sufficient, with no network access or external Python
packages required. Run without \texttt{-O}. Elapsed time and Python
version in the JSON are environment-dependent; the exact check counts
and rational outputs are deterministic.

\emph{Placement notes.} The program and its record are kept byte-identical
to the delivered files. Three facts about the placed layout matter.
\begin{enumerate}[label=\textup{(\roman*)},leftmargin=2em]
\item The program's only option is \texttt{-{}-output}. Its default is
\texttt{data/verification.json} in the directory that contains
\texttt{code/}, whatever the current directory. That file is the shipped
record, and the program writes its output there without checking whether
the file exists. The command in its docstring names the same path. A bare
\texttt{python code/verify.py} therefore overwrites the recorded evidence.
The elapsed-time field changes on every run, and the Python-version field
changes on another interpreter. Run the program on a copy of this
directory, or pass an explicit \texttt{-{}-output} path outside it.
\item The delivered \texttt{Makefile} is placed at \texttt{code/Makefile}.
Its recipes use paths relative to the report directory, so it is invoked
from there as \texttt{make -f code/Makefile verify}. That target writes
the new file \texttt{data/verification-rerun.json} and leaves the record
untouched.
\item The collection builds the article with
\texttt{latexmk -pdf -interaction=nonstopmode -halt-on-error article.tex}.
\end{enumerate}
At placement the program was rerun with Python~3.14.4, writing to a scratch
path. It reported \texttt{PASS} with 38,139 checks. Its output agreed with
\texttt{data/verification.json} in every field except the Python version
and the elapsed time.

\subsection{Current Lean coverage and the remaining route}
\label{epd:subsec:formalization}
A useful formalization boundary separates three layers. The first is
finite algebra: the divided-difference identity, simple-residue
uniqueness, exact valuation changes, and the orbit budget. The second
constructs the universal formal centers and evaluates them using the
finite-parameter Hahn map. The third proves the fiber decomposition,
the cylinder equivalence, and the topological dichotomy.

Once a general center-construction interface supplies
$F(x_s)=x_{\sh s}$ and the correct residues, the exact-fiber,
preperiodicity, and compact-subsystem proofs require no further
infinite summation. This isolates the support-heavy part of the remaining
Lean development.

The current mapping in \texttt{docs/FORMALIZATION.md} records
\cref{epd:def:orderunit,epd:lem:ideal} in
\nolinkurl{Surreal/HahnSeries/ScaleIdeal.lean}, over an arbitrary coefficient
field and ordered abelian group. It proves the convex scale subgroup,
the error ideal and its $q$-invariance, the intersection formula,
the order-unit criterion, and the coarsened valuation-ring description.
Clopenness is expressed by explicit valuation-ball conditions; no
Mathlib topology on Hahn series is presumed. The coarsening is represented
by a valuation subring, without an explicit ordered quotient isomorphism.

The positive-support lemma is also proved in the collection's
\nolinkurl{Surreal/HahnSeries/Neumann.lean}, including
\texttt{neumann\_positive} and \texttt{finite\_words\_of\_sum\_eq}.
The universal centers, exact fibers and later dynamical theorems are
not formalized by these modules. The report's written proofs and finite
regression checks are separate from this limited Lean coverage.

\section{Limitations and further research questions}
\label{epd:sec:questions}
The hypotheses in \eqref{epd:eq:polynomial} are doing real work. Simple
residue roots make the divided difference a unit, provide unique
branches, and make every derivative along the integral orbit have the
same valuation loss. A repeated residue root can require fractional
scales, ramification, or nonunique inverse branches. None of the
present theorems silently covers that situation.

The polynomial has finitely many coefficients and finitely many simple
branches. The finite-parameter evaluation proof does not authorize an
arbitrary countable family of independent perturbation parameters
without a new joint support hypothesis. Likewise, the general orbit
budget is for polynomials. Rational maps can repeatedly approach poles;
one cannot apply the polynomial bound to their denominators without
controlling those approaches.

\begin{question}[Nonuniform branch scales]\label{epd:q:branch-scales}
For a polynomial or rational map with finitely many inverse branches
whose exact gains are different positive values $\kappa_i$, can one
construct canonical itinerary centers and prove that the fiber of
$s$ is precisely the set of errors whose valuations exceed all partial
sums $\sum_{j<N}\kappa_{s_j}$? What finite-parameter support conditions
are sufficient when the branch equations are not simultaneously
normalized by one $q$?
\end{question}
The metric half of this proposed statement would follow from the same
pullback calculation. The nontrivial remaining issue is realization
of every admissible itinerary with a uniform support certificate.
The present article does not assume this generalization.

\begin{question}[A topology retaining both levels of information]\label{epd:q:enriched}
Can one give a useful enriched dynamical space that retains the compact
full-shift quotient and the nontrivial native infinitesimal fibers
simultaneously, without declaring the quotient topology to be the
wrong one? A precise comparison with higher-rank adic or other valued
geometric spaces would be preferable to relabeling the quotient as a
Berkovich space without a construction.
\end{question}

\begin{question}[Dynamics inside a single fiber]\label{epd:q:fiber}
For a periodic itinerary, what canonical normal forms classify the
restriction of the return map to $x_s+\Ik$? The valuation law is exact,
but it records only the size of a deviation, not its coefficients at
all finer scales. Which part of that finer structure survives the
flat deformations of \cref{epd:cor:flatdeform}?
\end{question}

These are proposed continuation problems arising from the proved
framework, not assertions that the same questions are established
open problems in the published literature.

\section{Conclusion}
The infinite word in an inverse-branch construction determines a point
only when the accumulated scales $n\kappa$ are cofinal in $\Gamma$, that is, when $\kappa$ is an order unit. Otherwise it determines
an exact affine infinitesimal fiber. Strong Hahn evaluation supplies
the canonical representative that topological convergence cannot
supply, and the exact divided-difference metric identifies every other
representative.

For the polynomial family studied here, this produces a sharp and
fully explicit separation between finite algebraic orbit data and
infinite native topological behavior. All preperiodicity equations
split in the original workspace. Enlarging the exponent group can
leave those equations and their roots unchanged while thickening
every itinerary fiber, destroying compact infinite invariant
subsystems, and turning the same center set from a Cantor space into
a discrete space. Infinitesimal coefficient changes beyond every
power of the dynamical scale preserve the symbolic partition but move
the actual fixed points.

The surreal example $F(x)=\omega(x^2-1)$ is therefore not merely an
ordinary real or complex dynamical system with a larger coefficient
field. Its behavior depends on which scales the chosen topology can
resolve. The correct conclusion is neither that ordinary chaos
survives verbatim nor that all dynamics becomes trivial: the symbolic
quotient, the exact fibers, the algebraic periodic locus, and the
ambient topology retain different, explicitly calculable information.

\appendix
\section{Dependency and claim ledger}
\label{epd:app:ledger}
\begin{table}[ht]
\centering\small
\caption{Logical dependencies and status of the principal results.}
\begin{tabularx}{\linewidth}{@{}p{0.30\linewidth}X@{}}
\toprule
Result & Dependencies and status\\\midrule
Positive-support lemma & Classical imported result; cited to Neumann. (Placement: proved in the collection's Lean library, \cref{epd:subsec:formalization}.)\\
Finite-parameter evaluation & Explicit consequence with coefficient-finiteness proof in \cref{epd:lem:evaluation}.\\
Universal centers & Formal implicit recursion, formal $Q$-adic completeness, and strong evaluation; proved in \cref{epd:thm:centers}.\\
Exact fibers & Centers plus the exact divided-difference metric; proposed new arbitrary-rank formulation, \cref{epd:thm:fibers}.\\
Preperiodicity splitting & Exact fibers and finite polynomial degree; proved in \cref{epd:thm:periodic}. Counting formulas are classical consequences.\\
Native topology and compact rigidity & Exact first mismatch, open error ideal, and elementary compactness; proposed new package, \cref{epd:thm:topology,epd:thm:compact}. (Placement: the proposal concerns set-sized workspaces; in $\No[\mathrm i]$ the set-level clauses are special cases of \texttt{a:thm:discrete}, \cref{epd:rem:discrete-special}.)\\
Scalar extension and transcendence & Universal support, preperiodicity splitting, and one-term valuation dominance; \cref{epd:thm:extension}.\\
Deformation and flat invisibility & Formal perturbation ideal before evaluation; \cref{epd:thm:deformation,epd:cor:flatdeform}.\\
Bounded orbit locus & Exact exterior valuation recurrence and order-unit criterion; \cref{epd:thm:bounded}.\\
General polynomial equicontinuity & Finite-degree growth budget; structural consequence of no order unit, \cref{epd:thm:general}.\\
Surreal realization & Classical Conway normal form plus the set-sized workspace proofs; \cref{epd:thm:surreal,epd:cor:surreal-all}.\\
\bottomrule
\end{tabularx}
\end{table}

All newly proved assertions are mathematical claims of this draft.
No independent peer review is claimed. The precise partial Lean coverage
is stated in \cref{epd:subsec:formalization}; it does not certify the
remaining dynamical theorems.
Finite tests support only the finite identities and regression cases
listed in \cref{epd:sec:verification}; they do not establish the general
statements in this ledger.

\section{A compact diagram of the topology change}
\label{epd:app:diagram}
\begin{center}
\begin{tikzpicture}[
 node distance=18mm and 17mm,
 every node/.style={align=center,font=\small},
 box/.style={draw=ink,rounded corners=2pt,fill=soft,
             text width=52mm,minimum height=18mm},
 >=Latex]
\node[box] (ground) {$\BI(F;K)=\{x_s\}$\\
$\kappa$ an order unit of $\Gamma$\\compact Cantor topology};
\node[box,right=of ground] (large) {$\BI(F;L)=\coprod_s(x_s+\Ik)$\\
$\kappa$ not an order unit of $\Delta$\\open fibers; centers discrete};
\node[box,below=of large] (shift) {$\Sigma_d$ with cylinder topology\\
full shift, entropy $\log d$};
\draw[->,thick,color=ink] (ground) -- node[above,text width=15mm]{field\\extension} (large);
\draw[->,thick,color=ink] (large) -- node[right,text width=37mm]{Hausdorff quotient\\of the scale uniformity} (shift);
\draw[->,thick,color=ink] (ground) |- node[below,pos=.72]{itinerary homeomorphism} (shift);
\end{tikzpicture}
\end{center}
The upper arrow is an algebraic inclusion of point sets; it is not
claimed continuous for the two intrinsic native topologies. The
right-hand arrow is continuous in the native topology but is not its
quotient map. The label on that arrow specifies the coarser uniformity
whose Hausdorff quotient actually is the full shift.

\section{Limitations and non-claims, collected}
\label{epd:app:nonclaims}
This appendix was added at placement. It gathers the limitations,
disclaimers, and non-claims that the manuscript states in several places
(its article, its delivered README, the delivered
\texttt{RESEARCH\_STATUS.md}, and the \texttt{limitations} field of
\texttt{data/verification.json}). Items S1--S28 come from the source and
were regrouped here. Their historical no-formalization wording is updated
to the current coverage in \cref{epd:subsec:formalization}.
Items P1--P3 were added at placement. Nothing in
the body is weakened by this list, and nothing outside it is claimed.

\subsection*{Status, priority, and provenance}
\begin{enumerate}[label=\textbf{S\arabic*.},leftmargin=2.8em]
\item The report is an unrefereed research draft. Its partial machine-checked
coverage is specified in \cref{epd:subsec:formalization}; the remaining
dynamical claims have no such certificate.
\item Priority is not certified. The fiber, compact-subsystem, and
extension/deformation package is \emph{proposed} as new after a targeted
literature search whose results were incomplete and sometimes noisy.
Failure to locate the statements is not proof that they have never
appeared, and the claims need independent expert review before
publication as new (\cref{epd:sec:intro}; \texttt{RESEARCH\_STATUS.md}).
\item No named published open conjecture is claimed solved. The three
questions of \cref{epd:sec:questions} are proposed continuations, not
established open problems.
\item Classical inputs are credited, not claimed. These are the
positive-support lemma, formal implicit-function recursion, valuation
coarsening, the one-sided full shift and its periodic counts and zeta
series, the rank-one Cantor construction, and Conway normal form
(\cref{epd:sec:intro,epd:sec:periodic,epd:sec:surreal}).
\item The delivered manuscript included no Lean file and ran no repository
build. The current collection does formalize the scale subgroup and
error-ideal lemma, as recorded in \cref{epd:subsec:formalization};
this is not a formalization of the entire article.
\item The repository audit was targeted: the inventories and the dynamics
report's description only, with no line-by-line review. No assertion of
complete absence of overlap across every file is made. During that original survey the repository was
neither rebuilt nor modified; its descriptions are provenance,
not validation (\cref{epd:subsec:repo-pin}; \texttt{RESEARCH\_STATUS.md};
bibliography).
\item The repository's Hahn support machinery and its distinctions among
analytic categories are not presented as discoveries
(\cref{epd:subsec:repo-pin}).
\item The article does not purport to resolve the repository's unrelated
birthday, automorphism, or differential-equation questions.
\end{enumerate}

\subsection*{Hypotheses}
\begin{enumerate}[label=\textbf{S\arabic*.},leftmargin=2.8em,resume]
\item The theorems need a degree-preserving reduction whose roots are
distinct and lie in $k$; for $k=\R$ they must be real. Repeated residue roots,
with fractional scales, ramification, or nonunique branches, are not
covered (\cref{epd:sec:questions}).
\item Only finitely many coefficients and branches are allowed. An
arbitrary countable family of independent perturbation parameters is not
covered without a new joint support hypothesis.
\item The orbit budget is for polynomials. Rational maps, which can
repeatedly approach poles, are not covered.
\item The nonuniform-branch-scale generalization of
\cref{epd:q:branch-scales} is not assumed anywhere.
\end{enumerate}

\subsection*{What the conclusions mean}
\begin{enumerate}[label=\textbf{S\arabic*.},leftmargin=2.8em,resume]
\item \emph{Valuation-expanding} does not mean repelling for a real-valued
absolute value. The expanding points are not identified with classical
Julia points (\cref{epd:sec:intro}, \cref{epd:rem:expansion}).
\item Integral nonescape and valuation-boundedness are different
properties (\cref{epd:sec:bounded}).
\item Strong Hahn summation is an algebraic support operation, not a
topological limit. Partial sums need not converge in the native topology,
and the nested approximants do not converge when $\kappa$ is not an order
unit (\cref{epd:rem:no-limits,epd:rem:nestedlimit,epd:rem:formal-ideal}).
\item The centers are canonical only relative to the fixed embedding
$k\hookrightarrow K$, the reduction $P_0$, and the universal construction.
No intrinsic uniqueness among set-theoretic sections of $\pi$ is asserted
(\cref{epd:sec:centers}).
\item In \cref{epd:thm:deformation}, no equality is asserted when the
leading numerator $D_\theta(c_{s_0})$ vanishes.
\item \Cref{epd:thm:compact} assumes $F(A)\subseteq A$, not $F(A)=A$, and
does not assert that every point is periodic. There is no conflict with
entropy monotonicity for compact factors, because the native space is
noncompact.
\item The full shift is the Hausdorff quotient of the coarsened
$\kappa$-scale uniformity, not the native quotient
(\cref{epd:prop:scalequotient}). In Appendix~\ref{epd:app:diagram} the extension
arrow is not claimed continuous, and the right-hand arrow is continuous but
not a quotient map.
\item $J_{\mathrm{aff},v}=\varnothing$ is a statement about the stated
affine native notion only. It is not claimed that Berkovich or spherical
Julia sets are empty, or that a rank-one Julia set disappears in its own
topology (\cref{epd:cor:empty-affine,epd:cor:surreal-all}).
\item No set of all surreal numbers is formed and no proper-class sum is
taken. Compactness in the surreal field refers to set-sized subsets. The
fine and workspace topologies are two topologies, not two incompatible
conclusions (\cref{epd:sec:surreal}).
\item Primes are formal $X$-derivatives with the coefficients held fixed.
No derivation on surreal coefficients is used, and no differential result
of Berarducci--Mantova is used (\cref{epd:sec:branches}, \cite{BM}).
\end{enumerate}

\subsection*{Computations}
\begin{enumerate}[label=\textbf{S\arabic*.},leftmargin=2.8em,resume]
\item The finite checks prove no infinite statement. They do not prove
infinite Hahn summability. The transfinite-support, topology,
compactness, and proper-class results are proved in the text, not checked
by the program. Nothing computed decides membership in $\Ik$
(\cref{epd:sec:verification}; \texttt{data/verification.json}).
\item The two center algorithms are different but share the defining
equation. Their agreement is not independent foundational verification.
\item The $O(N^3)$ operation count excludes coefficient bit growth.
\item The rank-two checks do not turn an approximate fixed-point
truncation into an exact Hahn fixed point.
\item The delivered finite verification uses no general surreal canonicalizer,
floating-point simulation, or proof-assistant certificate (\texttt{data/verification.json}).
\item Packaging (delivered README): the delivered PDF was not an accessible
tagged PDF, and no third-party books, research-paper PDFs, or
proof-assistant binaries were included. The PDF in this directory is
rebuilt from the source and is not tagged either.
\end{enumerate}

\subsection*{Added at placement}
\begin{enumerate}[label=\textbf{P\arabic*.},leftmargin=2.8em]
\item No result is transferred into or out of this report from
\texttt{dynamics-and-\allowbreak normal-forms}, \texttt{analysis},
\texttt{entire-functions-\allowbreak at-arbitrary-rank},
\texttt{hahn-tate-\allowbreak uniformization}, or \texttt{rank-one-\allowbreak berkovich}. This report
answers neither the Julia/Fatou exclusion nor the repelling-multiplier
exclusion of \texttt{dynamics-and-\allowbreak normal-forms}
(\cref{epd:subsec:collection}).
\item For the class $\No[\mathrm i]$ with its fine topology, the set-level
discreteness and compactness clauses are special cases of
\texttt{a:thm:discrete}. The proposal of novelty for those clauses concerns
set-sized workspaces (\cref{epd:rem:discrete-special}).
\item The order-unit dichotomy of \cref{epd:thm:topology} is not the
cofinality dichotomy \texttt{ent:\allowbreak cor:cofinality}. The convex
subgroup $H_\kappa$ shares its definition with $H_\alpha$ of
\texttt{hahn-tate-\allowbreak uniformization}. That identifies a
subgroup, not a theorem.
\end{enumerate}

\begin{thebibliography}{9}

\bibitem{Neumann}
B.~H. Neumann,
\emph{On ordered division rings},
Transactions of the American Mathematical Society \textbf{66} (1949),
202--252.
\href{https://doi.org/10.1090/S0002-9947-1949-0032593-5}{doi:10.1090/S0002-9947-1949-0032593-5}.
Used for the positive-support lemma for ordered series.

\bibitem{Gonshor}
H.~Gonshor,
\emph{An Introduction to the Theory of Surreal Numbers},
London Mathematical Society Lecture Note Series, vol.~110,
Cambridge University Press, 1986.
\href{https://www.cambridge.org/core/books/an-introduction-to-the-theory-of-surreal-numbers/312AE504A3E88E804054BFB390446374}{Publisher record}.
Classical background for Conway normal form and surreal arithmetic.

\bibitem{BM}
A.~Berarducci and V.~Mantova,
\emph{Surreal numbers, derivations and transseries},
Journal of the European Mathematical Society \textbf{20} (2018),
339--390.
\href{https://arxiv.org/abs/1503.00315}{arXiv:1503.00315}.
The freely available version explicitly describes the Hahn normal-form
structure; no differential result from that paper is used here.

\bibitem{Benedetto}
R.~L. Benedetto,
\emph{Non-Archimedean dynamics in dimension one: lecture notes},
Arizona Winter School, 2010, version of 9 March 2010.
\href{https://swc-math.github.io/aws/2010/2010BenedettoNotes-09Mar.pdf}{Freely available lecture notes}.
Definitions~4.15--4.16 and Example~4.39 provide the classical topological
and quadratic comparison, not the arbitrary-rank conclusions proved here.

\bibitem{Repo}
V.~Reshetnikov,
\emph{Surreal}, repository root and documentation inventories,
snapshot \texttt{a3124af79f66b8b9c196d76b4cbc5ac3938907c4},
consulted 22 September 2026.
\href{https://github.com/VladimirReshetnikov/Surreal/tree/a3124af79f66b8b9c196d76b4cbc5ac3938907c4}{Pinned repository};
\href{https://github.com/VladimirReshetnikov/Surreal/blob/a3124af79f66b8b9c196d76b4cbc5ac3938907c4/docs/README.md}{documentation inventory}.
Repository descriptions are provenance, not independent validation of
this article.

\bibitem{RepoDynamics}
\emph{Surcomplex Dynamics: Linearization Thresholds, Periods, and Normal Forms},
README for \texttt{docs/surcomplex/dynamics-and-normal-forms},
in the same pinned repository \cite{Repo}.
\href{https://github.com/VladimirReshetnikov/Surreal/blob/a3124af79f66b8b9c196d76b4cbc5ac3938907c4/docs/surcomplex/dynamics-and-normal-forms/README.md}{Pinned report description}.
Consulted for scope and overlap; the full merged manuscript was not
exhaustively audited.

\end{thebibliography}
\end{document}
